\chapter{Espaços métricos} \label{ms:chapter}

%%%%%%%%%%%%%%%%%%%%%%%%%%%%%%%%%%%%%%%%%%%%%%%%%%%%%%%%%%%%%%%%%%%%%%%%%%%%%%

\section{Espaços métricos}
\label{sec:metric}

\sectionnotes{1,5 aulas}

Como mencionamos na introdução, a ideia principal da análise é tomar
limites. Em \chapterref{seq:chapter}, aprendemos a calcular limites de
sequências de números reais. Em \chapterref{lim:chapter}, aprendemos a
calcular limites de funções quando uma variável real se aproxima de certo
número real.

Queremos calcular limites em contextos mais complicados. Por exemplo,
queremos considerar sequências de pontos no espaço tridimensional.
Queremos definir funções contínuas de várias variáveis.
Queremos até definir funções em espaços um pouco mais difíceis de descrever,
como a superfície da Terra. Ainda queremos falar de limites nesses espaços.

Por fim, vimos o limite de uma sequência de funções em
\chapterref{fs:chapter}.
Queremos unificar todas essas noções para não ter de demonstrar os mesmos
teoremas repetidas vezes em cada contexto. O conceito de espaço métrico é
uma ferramenta elementar, mas poderosa, em análise. Embora não seja suficiente
para descrever todos os tipos de limite encontrados na análise moderna, ele
nos permite avançar bastante.

\begin{defn}
Seja $X$ um conjunto, e seja
$d \colon X \times X \to \R$
uma função tal que, para quaisquer $x,y,z \in X$,
\begin{enumerate}[(i)]
%
\item \label{metric:pos}
%mbxSTARTIGNORE
\makebox[2.8in][l]{$d(x,y) \geq 0$.}
%mbxENDIGNORE
%mbxlatex $d(x,y) \geq 0$. \qquad
(\emph{não negatividade})\index{não negatividade!de uma métrica}
%
\item \label{metric:zero}
%mbxSTARTIGNORE
\makebox[2.8in][l]{$d(x,y) = 0$ se, e somente se, $x = y$.}
%mbxENDIGNORE
%mbxlatex $d(x,y) = 0$ se, e somente se, $x = y$. \qquad
(\emph{\myindex{identidade dos indiscerníveis}})
%
\item \label{metric:com}
%mbxSTARTIGNORE
\makebox[2.8in][l]{$d(x,y) = d(y,x)$.} 
%mbxENDIGNORE
%mbxlatex $d(x,y) = d(y,x)$. \qquad
(\emph{simetria})\index{simetria!de uma métrica}
%
\item \label{metric:triang}
%mbxSTARTIGNORE
\makebox[2.8in][l]{$d(x,z) \leq d(x,y)+ d(y,z)$.}
%mbxENDIGNORE
%mbxlatex $d(x,z) \leq d(x,y)+ d(y,z)$. \qquad
(\emph{\myindex{desigualdade triangular}})
\end{enumerate}
O par $(X,d)$ é chamado de \emph{\myindex{espaço métrico}}. A função $d$ é
chamada de \emph{\myindex{métrica}} ou de \emph{\myindex{função distância}}.
Às vezes, escrevemos simplesmente $X$ para indicar o espaço métrico, em vez
de $(X,d)$, quando a métrica fica clara pelo contexto.
\end{defn}

A ideia geométrica é que $d$ representa a distância entre dois pontos.
Os itens \ref{metric:pos}--\ref{metric:com} têm uma interpretação geométrica
óbvia: a distância é sempre não negativa; o único ponto cuja distância a
$x$ é zero é o próprio $x$; e a distância de $x$ a $y$ é igual à distância
de $y$ a $x$. A desigualdade triangular \ref{metric:triang} tem a
interpretação ilustrada na \figureref{fig:mstriang}.
\begin{myfigureht}
\myincludepdft{ms-triang}{%
Um triângulo com vértices x, y e z.
A distância entre x e z é indicada por d(x,z); a distância de x a y,
por d(x,y); e a distância de y a z, por d(y,z).
Uma seta de x a z está identificada como "mais curto", e uma seta que vai
de x a y e depois segue até z está identificada como "mais longo".}
\caption{Diagrama da desigualdade triangular em espaços métricos.\label{fig:mstriang}}
\end{myfigureht}

Para desenhar, é conveniente representar figuras e diagramas no plano, usando
a distância euclidiana como métrica.
No entanto, esse é apenas um espaço métrico particular. O fato de algo parecer
claro ao observar um desenho não significa que seja verdadeiro em todo espaço
métrico.
Talvez a intuição da geometria euclidiana esteja desviando sua atenção, embora
o conceito de espaço métrico seja muito mais geral.

Vejamos alguns exemplos de espaços métricos.

\begin{example}
O conjunto dos números reais $\R$ é um espaço métrico com a métrica
\begin{equation*}
d(x,y) \coloneqq \sabs{x-y} .
\end{equation*}
É fácil verificar os itens \ref{metric:pos}--\ref{metric:com} da definição.
A desigualdade triangular \ref{metric:triang} decorre imediatamente da
desigualdade triangular usual para números reais:
\begin{equation*}
d(x,z) = \sabs{x-z} = 
\sabs{x-y+y-z} \leq
\sabs{x-y}+\sabs{y-z} =
d(x,y)+ d(y,z) .
\end{equation*}
Essa métrica é a \emph{\myindex{métrica usual em $\R$}}. Se falarmos de $\R$
como espaço métrico sem especificar uma métrica, estaremos nos referindo a
essa métrica.
\end{example}

\begin{example}
Também podemos definir outra métrica no conjunto dos números reais.
Por exemplo, consideremos o conjunto dos números reais $\R$ com a métrica
\begin{equation*}
d(x,y) \coloneqq
\frac{\sabs{x-y}}{\sabs{x-y}+1} .
\end{equation*}
Novamente, é fácil verificar os itens \ref{metric:pos}--\ref{metric:com}.
A desigualdade triangular \ref{metric:triang} é um pouco mais difícil.
Observe que $d(x,y) = \varphi\bigl(\sabs{x-y}\bigr)$, em que
$\varphi(t) = \frac{t}{t+1}$, e que $\varphi$ é uma função crescente (sua
derivada é positiva; veja a \figureref{fig:tovertp1}). Portanto,
\begin{equation*}
\begin{split}
d(x,z) = \varphi\bigl(\sabs{x-z}\bigr)
& = 
\varphi\bigl(\sabs{x-y+y-z}\bigr) \\
& \leq
\varphi\bigl(\sabs{x-y}+\sabs{y-z}\bigr)
\\
& =
\frac{\sabs{x-y}+\sabs{y-z}}{\sabs{x-y}+\sabs{y-z}+1} \\
& =
\frac{\sabs{x-y}}{\sabs{x-y}+\sabs{y-z}+1} +
\frac{\sabs{y-z}}{\sabs{x-y}+\sabs{y-z}+1}
\\
& \leq
\frac{\sabs{x-y}}{\sabs{x-y}+1} +
\frac{\sabs{y-z}}{\sabs{y-z}+1} =
d(x,y)+ d(y,z) .
\end{split}
\end{equation*}
Assim, $d$ é uma métrica e fornece um exemplo de métrica não usual em $\R$.
Nessa métrica, $d(x,y) < 1$ para todos $x,y \in \R$. Isto é, quaisquer dois
pontos estão a uma distância menor que 1 unidade.
\begin{myfigureht}
\myincludegraphics{tovertp1graph}{%
Gráfico de uma função que parte da origem subindo diagonalmente e depois se
curva lentamente, tornando-se quase horizontal na extremidade direita.
Uma linha tracejada acima do gráfico indica uma assíntota horizontal.}
\caption{Gráfico de $\frac{t}{t+1}$ para $t$ positivo, com assíntota em 1.\label{fig:tovertp1}}
\end{myfigureht}
\end{example}

Um espaço métrico importante é o
\emph{\myindex{espaço euclidiano}} $n$-dimensional
\glsadd{not:euclidspace}
$\R^n = \R \times \R \times \cdots \times \R$. Usamos a seguinte notação
para os pontos: $x =(x_1,x_2,\ldots,x_n) \in \R^n$. Não escreveremos
$\vec{x}$ nem $\mathbf{x}$ para indicar um ponto de $\R^n$, como é comum no
cálculo de várias variáveis; simplesmente lhe daremos um nome, como $x$, e
lembraremos que $x$ é um elemento de $\R^n$.
Também escreveremos simplesmente $0 \in \R^n$ para indicar o ponto
$(0,0,\ldots,0)$. Antes de tornar $\R^n$ um espaço métrico, demonstraremos
uma desigualdade importante: a desigualdade de Cauchy--Schwarz.

\begin{lemma}[\myindex{Cauchy--Schwarz inequality}%
\footnote{%
Às vezes, ela é chamada de \myindex{desigualdade de Cauchy--Bunyakovsky--Schwarz}.
\href{https://en.wikipedia.org/wiki/Hermann_Schwarz}{Karl Hermann Amandus
Schwarz} (1843--1921) foi um matemático alemão e
\href{https://en.wikipedia.org/wiki/Viktor_Bunyakovsky}{Viktor Yakovlevich
Bunyakovsky} (1804--1889) foi um matemático ucraniano. O que enunciamos
deveria, a rigor, ser chamado de desigualdade de Cauchy, pois Bunyakovsky e
Schwarz demonstraram versões em dimensão infinita.}]
Sejam $x =(x_1,x_2,\ldots,x_n) \in \R^n$ e $y =(y_1,y_2,\ldots,y_n) \in \R^n$.
Então,
\begin{equation*}
{\biggl( \sum_{k=1}^n x_k y_k \biggr)}^2
\leq
\biggl(\sum_{k=1}^n x_k^2 \biggr)
\biggl(\sum_{k=1}^n y_k^2 \biggr) .
\end{equation*}
\end{lemma}

\begin{proof}
O quadrado de um número real é não negativo. Portanto, uma soma de quadrados
é não negativa:
\begin{equation*}
\begin{split}
0 & \leq 
\sum_{k=1}^n \sum_{\ell=1}^n {(x_k y_\ell - x_\ell y_k)}^2
\\
& =
\sum_{k=1}^n \sum_{\ell=1}^n \bigl( x_k^2 y_\ell^2 + x_\ell^2 y_k^2 - 2 x_k
x_\ell y_k
y_\ell \bigr)
\\
& =
\biggl( \sum_{k=1}^n x_k^2 \biggr)
\biggl( \sum_{\ell=1}^n y_\ell^2 \biggr)
+
\biggl( \sum_{k=1}^n y_k^2 \biggr)
\biggl( \sum_{\ell=1}^n x_\ell^2 \biggr)
-
2
\biggl( \sum_{k=1}^n x_k y_k \biggr)
\biggl( \sum_{\ell=1}^n x_\ell y_\ell \biggr) .
\end{split}
\end{equation*}
Renomeando os índices e dividindo por 2, obtemos exatamente o que queríamos:
\begin{equation*}
0 \leq 
\biggl( \sum_{k=1}^n x_k^2 \biggr)
\biggl( \sum_{k=1}^n y_k^2 \biggr)
-
{\biggl( \sum_{k=1}^n x_k y_k \biggr)}^2 .
\qedhere
\end{equation*}
\end{proof}

\begin{example}
Vamos definir a métrica usual\index{métrica usual em $\R^n$} em $\R^n$:
\begin{equation*}
d(x,y) \coloneqq
\sqrt{
{(x_1-y_1)}^2 + 
{(x_2-y_2)}^2 + 
\cdots +
{(x_n-y_n)}^2
} =
\sqrt{
\sum_{k=1}^n
{(x_k-y_k)}^2 
} .
\end{equation*}
Para $n=1$, isto é, na reta real, essa métrica coincide com a que definimos
acima. Para $n > 1$, como antes, o único ponto delicado da definição a
verificar é a desigualdade triangular.
É mais simples trabalhar com o quadrado da métrica. Na estimativa a seguir,
observe o uso da desigualdade de Cauchy--Schwarz.
\begin{equation*}
\begin{split}
{\bigl(d(x,z)\bigr)}^2 & =
\sum_{k=1}^n
{(x_k-z_k)}^2 
\\
& =
\sum_{k=1}^n
{(x_k-y_k+y_k-z_k)}^2 
\\
& =
\sum_{k=1}^n
\Bigl(
{(x_k-y_k)}^2+{(y_k-z_k)}^2 + 2(x_k-y_k)(y_k-z_k)
\Bigr)
\\
& =
\sum_{k=1}^n
{(x_k-y_k)}^2
+
\sum_{k=1}^n
{(y_k-z_k)}^2 
+
2
\sum_{k=1}^n
(x_k-y_k)(y_k-z_k)
\\
& \leq
\sum_{k=1}^n
{(x_k-y_k)}^2
+
\sum_{k=1}^n
{(y_k-z_k)}^2 
+
2
\sqrt{
\sum_{k=1}^n
{(x_k-y_k)}^2
\sum_{k=1}^n
{(y_k-z_k)}^2
}
\\
& =
{\left(
\sqrt{
\sum_{k=1}^n
{(x_k-y_k)}^2
}
+
\sqrt{
\sum_{k=1}^n
{(y_k-z_k)}^2 
}
\right)}^2
=
{\bigl( d(x,y) + d(y,z) \bigr)}^2 .
\end{split}
\end{equation*}
Como a raiz quadrada é uma função crescente, a desigualdade é preservada ao
extrairmos a raiz quadrada dos dois lados, e obtemos a desigualdade triangular.
\end{example}

\begin{example}
O conjunto dos números complexos $\C$ é formado pelos números $z = x+iy$,
em que $x$ e $y$ pertencem a $\R$. Ao impormos $i^2 = -1$, transformamos
$\C$ em um corpo.
Para calcular limites, consideramos $\C$ como o espaço métrico $\R^2$, em
que $z=x+iy \in \C$ corresponde a $(x,y) \in \R^2$.
Para $z=x+iy$, definimos o \emph{\myindex{módulo complexo}}\index{módulo}
por $\sabs{z} \coloneqq \sqrt{x^2+y^2}$.
Então, para dois números complexos $z_1 = x_1 + iy_1$ e
$z_2 = x_2 + iy_2$, a distância é
\begin{equation*}
d(z_1,z_2) = \sqrt{{(x_1-x_2)}^2+ {(y_1-y_2)}^2} = \sabs{z_1-z_2}.
\end{equation*}
Além disso, ao trabalhar com números complexos, muitas vezes é conveniente
escrever a métrica em termos do \emph{\myindex{conjugado complexo}}: o
conjugado de $z=x+iy$ é $\bar{z} \coloneqq x-iy$. Então,
${\sabs{z}}^2 = x^2 +y^2 = z\bar{z}$, e, portanto, ${\sabs{z_1-z_2}}^2 =
(z_1-z_2)\overline{(z_1-z_2)}$.
\end{example}

\begin{example}
Um exemplo útil para ter em mente é a \emph{\myindex{métrica discreta}}.
Para qualquer conjunto $X$, definimos
\begin{equation*}
d(x,y) \coloneqq
\begin{cases}
1 & \text{se } x \neq y, \\
0 & \text{se } x = y.
\end{cases}
\end{equation*}
Isto é, quaisquer dois pontos distintos estão à mesma distância. Quando $X$
é finito, podemos desenhar um diagrama, como o da \figureref{fig:msdiscmetric}.
É claro que, no diagrama, as distâncias não são as distâncias euclidianas
usuais do plano.
As coisas ficam mais sutis quando $X$ é infinito, como no caso dos números
reais.
\begin{myfigureht}
\myincludepdft{msdiscmetric}{%
Diagrama com cinco pontos a, b, c, d e e. Cada ponto está ligado a
todos os demais por um segmento, e cada segmento está identificado com o
número 1.}
\caption{Exemplo de espaço métrico discreto $\{ a,b,c,d,e \}$: a distância
entre quaisquer dois pontos distintos é $1$.\label{fig:msdiscmetric}}
\end{myfigureht}

Embora esse exemplo talvez raramente apareça na prática, ele fornece um
\myquote{teste de plausibilidade} útil. Se você fizer uma afirmação sobre
espaços métricos, experimente verificá-la para a métrica discreta.
Demonstrar que $(X,d)$ é de fato um espaço métrico fica como exercício.
\end{example}

\begin{example} \label{example:msC01}
Seja $C\bigl([a,b],\R\bigr)$\glsadd{not:contfuncs} o conjunto das funções
contínuas de valores reais definidas no intervalo $[a,b]$.
Definimos a métrica em $C\bigl([a,b],\R\bigr)$ por
\begin{equation*}
d(f,g) \coloneqq \sup_{x \in [a,b]} \babs{f(x)-g(x)} .
\end{equation*}
Verifiquemos as propriedades. Primeiro, $d(f,g)$ é finito, pois
$\babs{f(x)-g(x)}$ é uma função contínua em um intervalo fechado e limitado
$[a,b]$ e, portanto, é limitada.
É claro que $d(f,g) \geq 0$, pois se trata do supremo de números não negativos.
Se $f = g$, então $\babs{f(x)-g(x)} = 0$ para todo $x$ e, logo,
$d(f,g) = 0$. Reciprocamente, se $d(f,g) = 0$, então, para todo $x$,
$\babs{f(x)-g(x)} \leq d(f,g) = 0$; portanto, $f(x) = g(x)$ para todo $x$,
e assim $f=g$. Também é trivial que $d(f,g) = d(g,f)$. Para demonstrar a
desigualdade triangular, usamos a desigualdade triangular usual.
\begin{equation*}
\begin{split}
d(f,g) & =
\sup_{x \in [a,b]} \babs{f(x)-g(x)} =
\sup_{x \in [a,b]} \babs{f(x)-h(x)+h(x)-g(x)}
\\
& \leq
\sup_{x \in [a,b]} \Bigl( \babs{f(x)-h(x)}+\babs{h(x)-g(x)} \Bigr)
\\
& \leq
\sup_{x \in [a,b]} \babs{f(x)-h(x)}+
\sup_{x \in [a,b]} \babs{h(x)-g(x)} = d(f,h) + d(h,g) .
\end{split}
\end{equation*}
Quando consideramos $C\bigl([a,b],\R\bigr)$ como espaço métrico sem mencionar
uma métrica específica, estamos nos referindo a essa métrica.
Observe que $d(f,g) = \snorm{f-g}_{[a,b]}$, a norma uniforme da
\defnref{def:unifnorm}.

À primeira vista, esse exemplo pode parecer esotérico, mas trabalhar com
espaços como $C\bigl([a,b],\R\bigr)$ está no cerne de boa parte da análise
moderna. Considerar conjuntos de funções como espaços métricos nos permite
abstrair muitos detalhes técnicos e demonstrar com menos trabalho resultados
poderosos, como o \hyperref[thm:fs:picard]{teorema de Picard}.
\end{example}

\begin{example}
Outro exemplo útil de espaço métrico é a esfera com uma métrica geralmente
chamada de \emph{\myindex{distância ortodrômica}}.
Seja $S^2$ a \myindex{esfera unitária}\index{esfera} em $\R^3$, isto é,
$S^2 \coloneqq \{ x \in \R^3 : x_1^2+x_2^2+x_3^2 = 1 \}$.
Tomemos $x$ e $y$ em $S^2$, tracemos uma reta pela origem e por $x$ e outra
reta pela origem e por $y$, e seja $\theta$ o ângulo entre essas retas.
Definimos então $d(x,y) \coloneqq \theta$. Veja a
\figureref{fig:spheremetric}.
A lei dos cossenos da geometria vetorial fornece
$d(x,y) = \arccos(x_1y_1+x_2y_2+x_3y_3)$.
É relativamente fácil ver que essa função satisfaz as três primeiras
propriedades de uma métrica.
A desigualdade triangular é mais difícil de demonstrar e exige um pouco mais
de trigonometria e álgebra linear do que pretendemos usar agora; por isso,
deixaremos a demonstração de lado.

\begin{myfigureht}
\myincludepdft{spheremetric}{%
Diagrama de uma esfera no espaço tridimensional, identificada por S ao
quadrado. A origem está indicada por 0, e os pontos x e y pertencem à esfera.
Há uma reta entre a origem e x, outra entre a origem e y, e o ângulo entre
elas está identificado por teta.}
\caption{Distância ortodrômica na esfera
unitária.\label{fig:spheremetric}}
\end{myfigureht}

Essa distância é o menor caminho entre pontos da esfera quando só podemos
percorrer a própria esfera. É fácil generalizá-la para esferas de qualquer
diâmetro. Se considerarmos uma esfera de raio $r$, definimos a distância por
$d(x,y) \coloneqq r \theta$. Por exemplo, essa é a distância usual para
calcular distâncias sobre a superfície da Terra, como a distância percorrida
por um avião entre Londres e Los Angeles.
\end{example}

Muitas vezes, é útil considerar um subconjunto de um espaço métrico maior
como um espaço métrico por si só. Obtemos a proposição a seguir, cuja
demonstração é imediata.

\begin{prop}
Seja $(X,d)$ um espaço métrico e $Y \subset X$. Então, a restrição
$d|_{Y \times Y}$ é uma métrica em $Y$.
\end{prop}

\begin{defn}
Se $(X,d)$ é um espaço métrico, $Y \subset X$ e $d' \coloneqq d|_{Y \times Y}$,
então $(Y,d')$ é chamado de \emph{\myindex{subespaço}} de $(X,d)$.
\end{defn}

É comum escrever simplesmente $d$ para a métrica em $Y$, pois ela é a
restrição da métrica em $X$. Às vezes, dizemos que $d'$ é a
\emph{\myindex{métrica de subespaço}} e que $Y$ tem a
\emph{\myindex{topologia de subespaço}}.

\medskip

Um subconjunto dos números reais é limitado quando todos os seus elementos
distam de 0 no máximo uma constante fixa.
Em um espaço métrico qualquer, talvez não haja um ponto 0 natural e fixo,
mas isso não importa para a definição de limitação.

\begin{defn}
Seja $(X,d)$ um espaço métrico. Dizemos que um subconjunto $S \subset X$ é
\emph{limitado}\index{conjunto limitado} se existem $p \in X$ e $B \in \R$
tais que
\begin{equation*}
d(p,x) \leq B \quad \text{para todo } x \in S.
\end{equation*}
Dizemos que $(X,d)$ é limitado se o próprio $X$ é um subconjunto limitado.
\end{defn}

Por exemplo, o conjunto dos números reais com a métrica usual não é um
espaço métrico limitado. Não é difícil ver que um subconjunto dos números
reais é limitado no sentido de \chapterref{rn:chapter} se, e somente se,
é limitado como subconjunto do espaço métrico dos números reais com a métrica
usual.

Por outro lado, se considerarmos os números reais com a métrica discreta,
obtemos um espaço métrico limitado. Na verdade, qualquer conjunto com a
métrica discreta é limitado.

Há outras maneiras equivalentes de generalizar a noção de conjunto limitado,
que ficam como exercícios. Suponha que $X$ não seja vazio, para evitar uma
questão técnica. Então, $S \subset X$ ser limitado equivale a qualquer uma
das condições a seguir:
\begin{enumerate}[(i)]
\item
Para todo $p \in X$, existe $B > 0$ tal que $d(p,x) \leq B$ para todo $x \in S$.
\item
\glsadd{not:diam}%
$\operatorname{diam}(S) \coloneqq \sup \bigl\{ d(x,y) : x,y \in S \bigr\} < \infty$.
\end{enumerate}
A quantidade $\operatorname{diam}(S)$ é chamada de
\emph{\myindex{diâmetro}} de um conjunto e, em geral, só é definida para
conjuntos não vazios.

\subsection{Exercícios}

\begin{exercise}
Demonstre que, para todo conjunto $X$, a métrica discreta ($d(x,y) = 1$ se
$x \neq y$ e $d(x,x) = 0$) de fato define um espaço métrico $(X,d)$.
\end{exercise}

\begin{exercise}
Seja $X \coloneqq \{ 0 \}$ um conjunto. É possível torná-lo um espaço métrico?
\end{exercise}

\begin{exercise}
Seja $X \coloneqq \{ a, b \}$ um conjunto. É possível torná-lo dois espaços
métricos distintos? (Defina duas métricas distintas sobre esse conjunto.)
\end{exercise}

\begin{exercise}
Seja o conjunto $X \coloneqq \{ A, B, C \}$ o conjunto de três edifícios do
campus. Suponha que queiramos definir a distância como o tempo necessário
para caminhar de um edifício a outro.
O percurso entre os edifícios $A$ e $B$ leva 5 minutos em qualquer direção.
No entanto, o edifício $C$ fica em uma colina, e de $A$ são necessários
10 minutos e de $B$, 15 minutos, para chegar até $C$. Por outro lado,
de $C$ até $A$ são 5 minutos e de $C$ até $B$ são 7 minutos, pois o caminho é
descida. Essas distâncias definem uma métrica? Se sim, demonstre; caso
contrário, explique por quê.
\end{exercise}

\begin{exercise}
Suponha que $(X,d)$ seja um espaço métrico e que
$\varphi \colon [0,\infty) \to \R$ seja uma função crescente tal que
$\varphi(t) \geq 0$ para todo $t$ e $\varphi(t) = 0$ se, e somente se,
$t=0$. Suponha também que $\varphi$ seja \emph{\myindex{subaditiva}}, isto é,
$\varphi(s+t) \leq \varphi(s)+\varphi(t)$.
Mostre que, com $d'(x,y) \coloneqq \varphi\bigl(d(x,y)\bigr)$, obtemos um
novo espaço métrico $(X,d')$.
\end{exercise}

\begin{exercise} \label{exercise:mscross}
Sejam $(X,d_X)$ e $(Y,d_Y)$ espaços métricos.
\begin{enumerate}[a)]
\item
Mostre que $(X \times Y,d)$, com
$d\bigl( (x_1,y_1), (x_2,y_2) \bigr) \coloneqq d_X(x_1,x_2) + d_Y(y_1,y_2)$
é um espaço métrico.
\item
Mostre que $(X \times Y,d)$, com
$d\bigl( (x_1,y_1), (x_2,y_2) \bigr) \coloneqq \max \bigl\{ d_X(x_1,x_2) ,
d_Y(y_1,y_2) \bigr\}$
é um espaço métrico.
\end{enumerate}
\end{exercise}

\begin{exercise}
Seja $X$ o conjunto das funções contínuas em $[0,1]$. Seja $\varphi \colon
[0,1] \to (0,\infty)$ uma função contínua. Defina
\begin{equation*}
d(f,g) \coloneqq \int_0^1 \babs{f(x)-g(x)}\varphi(x)\,dx .
\end{equation*}
Demonstre que $(X,d)$ é um espaço métrico.
\end{exercise}

\begin{exercise} \label{exercise:mshausdorffpseudo}
Seja $(X,d)$ um espaço métrico. Para subconjuntos limitados não vazios $A$ e
$B$, defina
\begin{equation*}
d(x,B) \coloneqq \inf \bigl\{ d(x,b) : b \in B \bigr\}
\qquad \text{e} \qquad
d(A,B) \coloneqq \sup \bigl\{ d(a,B) : a \in A \bigr\} .
\end{equation*}
Defina então a \emph{\myindex{métrica de Hausdorff}} por
\begin{equation*}
d_H(A,B) \coloneqq \max \bigl\{ d(A,B) , d(B,A) \bigr\} .
\end{equation*}
Observação: podemos definir $d_H$ para quaisquer subconjuntos não vazios se
admitirmos os reais estendidos.
\begin{enumerate}[a)]
\item
Seja $Y \subset \sP(X)$ o conjunto dos subconjuntos limitados não vazios.
Demonstre que $(Y,d_H)$ é o que chamamos de
\emph{\myindex{espaço pseudométrico}}: $d_H$ satisfaz as propriedades métricas
\ref{metric:pos},
\ref{metric:com}, 
\ref{metric:triang}, e também
$d_H(A,A) = 0$ para todo $A \in Y$.
\item
Mostre por meio de um exemplo que a própria função $d$ não é simétrica, isto
é, $d(A,B) \neq d(B,A)$.
\item
Encontre um espaço métrico $X$ e dois subconjuntos limitados não vazios e
distintos $A$ e $B$ tais que $d_H(A,B) = 0$.
\end{enumerate}
\end{exercise}

\begin{exercise}
Seja $(X,d)$ um espaço métrico não vazio e $S \subset X$ um subconjunto.
Demonstre:
\begin{enumerate}[a)]
\item
$S$ é limitado se, e somente se, para todo $p \in X$ existe $B > 0$ tal que
$d(p,x) \leq B$ para todo $x \in S$.
\item
Um subconjunto $S$ não vazio é limitado se, e somente se,
$\operatorname{diam}(S) \coloneqq \sup \{ d(x,y) : x,y \in S \} < \infty$.
\end{enumerate}
\end{exercise}

\begin{exercise}
\pagebreak[2]
\leavevmode
\begin{enumerate}[a)]
\item
Trabalhando em $\R$, calcule $\operatorname{diam}\bigl([a,b]\bigr)$.
\item
Trabalhando em $\R^n$, para todo $r > 0$, seja
$B_r \coloneqq \{ x_1^2+x_2^2+\cdots+x_n^2 < r^2 \}$. Calcule
$\operatorname{diam}(B_r)$.
\item
Suponha que $(X,d)$ seja um espaço métrico com pelo menos dois pontos, que
$d$ seja a métrica discreta e que $p \in X$.
Calcule $\operatorname{diam}(\{ p \})$ e $\operatorname{diam}(X)$ e, em
seguida, conclua que $(X,d)$ é limitado.
\end{enumerate}
\end{exercise}

\begin{exercise}
\pagebreak[2]
\leavevmode
\begin{enumerate}[a)]
\item
Encontre uma métrica $d$ em $\N$ tal que $\N$ seja um conjunto ilimitado em
$(\N,d)$.
\item
Encontre uma métrica $d$ em $\N$ tal que $\N$ seja um conjunto limitado em
$(\N,d)$.
\item
Encontre uma métrica $d$ em $\N$ tal que, para todo $n \in \N$ e todo
$\epsilon > 0$, exista $m \in \N$, $m \neq n$, tal que $d(n,m) < \epsilon$.
\end{enumerate}
\end{exercise}

\begin{exercise} \label{exercise:C1ab}
Seja $C^1\bigl([a,b],\R\bigr)$\glsadd{not:contdifffuncs} o conjunto das
funções continuamente diferenciáveis em $[a,b]$.
Defina
\begin{equation*}
d(f,g) \coloneqq \snorm{f-g}_{[a,b]} + \snorm{f'-g'}_{[a,b]},
\end{equation*}
em que $\snorm{\cdot}_{[a,b]}$ é a norma uniforme. Demonstre que $d$ é uma
métrica.
\end{exercise}

\begin{samepage}
\begin{exercise}
Considere $\ell^2$, o conjunto das sequências $\{ x_n \}_{n=1}^\infty$ de
números reais tais que $\sum_{n=1}^\infty x_n^2 < \infty$.
\begin{enumerate}[a)]
\item
Demonstre a \myindex{desigualdade de Cauchy--Schwarz} para duas sequências
$\{x_n \}_{n=1}^\infty$ e $\{ y_n \}_{n=1}^\infty$ em $\ell^2$: demonstre
que $\sum_{n=1}^\infty x_n y_n$ converge (absolutamente) e que
\begin{equation*}
{\biggl( \sum_{n=1}^\infty x_n y_n \biggr)}^2
\leq
\biggl(\sum_{n=1}^\infty x_n^2 \biggr)
\biggl(\sum_{n=1}^\infty y_n^2 \biggr) .
\end{equation*}
\item
Demonstre que $\ell^2$ é um espaço métrico com a métrica
$d(x,y) \coloneqq \sqrt{\sum_{n=1}^\infty {(x_n-y_n)}^2}$.
Dica: não se esqueça de mostrar que a série que define $d(x,y)$ sempre
converge para um número finito.
\end{enumerate}
\end{exercise}
\end{samepage}

%%%%%%%%%%%%%%%%%%%%%%%%%%%%%%%%%%%%%%%%%%%%%%%%%%%%%%%%%%%%%%%%%%%%%%%%%%%%%%

\sectionnewpage
\section{Conjuntos abertos e fechados}
\label{sec:mettop}

%mbxINTROSUBSECTION

\sectionnotes{2 aulas}

\subsection{Topologia}

Antes de tratar da convergência,
definiremos a \emph{\myindex{topologia}}.
Isto é,
definiremos as noções de conjunto aberto e conjunto fechado.
Antes disso,
definiremos dois conjuntos abertos e fechados especiais.

\begin{defn}
Sejam $(X,d)$ um espaço métrico, $x \in X$ e $\delta > 0$. Definimos a
\emph{\myindex{bola aberta}}, ou simplesmente \emph{\myindex{bola}}, de raio
$\delta$ centrada em $x$ por
\glsadd{not:openball}
\begin{equation*}
B(x,\delta) \coloneqq \bigl\{ y \in X : d(x,y) < \delta \bigr\} .
\end{equation*}
Definimos a \emph{\myindex{bola fechada}} por
\glsadd{not:closedball}
\begin{equation*}
C(x,\delta) \coloneqq \bigl\{ y \in X : d(x,y) \leq \delta \bigr\} .
\end{equation*}
\end{defn}

Ao trabalhar com espaços métricos diferentes, às vezes é essencial destacar
em qual deles está a bola. Fazemos isso escrevendo
$B_X(x,\delta) \coloneqq B(x,\delta)$ ou
$C_X(x,\delta) \coloneqq C(x,\delta)$.

\begin{example}
Consideremos o espaço métrico $\R$ com a métrica usual. Para
$x \in \R$ e $\delta > 0$,
\begin{equation*}
B(x,\delta) = (x-\delta,x+\delta) \qquad \text{e} \qquad
C(x,\delta) = [x-\delta,x+\delta] .
\end{equation*}
\end{example}

\begin{example}
Tenha cuidado ao trabalhar em um subespaço. Consideremos o espaço métrico
$[0,1]$ como subespaço de $\R$. Então, em $[0,1]$,
\begin{equation*}
B(0,\nicefrac{1}{2}) = B_{[0,1]}(0,\nicefrac{1}{2})
= \bigl\{ y \in [0,1] : \sabs{0-y} < \nicefrac{1}{2} \bigr\}
= [0,\nicefrac{1}{2}) .
\end{equation*}
Isso é diferente de $B_{\R}(0,\nicefrac{1}{2}) =
(\nicefrac{-1}{2},\nicefrac{1}{2})$.
O importante é ter em mente em qual espaço métrico estamos trabalhando.
\end{example}

\begin{defn}
Seja $(X,d)$ um espaço métrico. Um subconjunto $V \subset X$ é
\emph{aberto}\index{conjunto aberto} se, para todo $x \in V$, existe
$\delta > 0$ tal que $B(x,\delta) \subset V$. Veja a
\figureref{fig:msopenset}. Um subconjunto $E \subset X$ é
\emph{fechado}\index{conjunto fechado} se seu complementar
$E^c = X \setminus E$ é aberto.
Quando o espaço ambiente $X$ não está claro pelo contexto,
dizemos que \emph{$V$ é aberto em $X$} e que \emph{$E$ é fechado em $X$}.

Se $x \in V$ e $V$ é aberto, dizemos que $V$ é uma
\emph{\myindex{vizinhança aberta}} de $x$ (ou, às vezes, simplesmente uma
\emph{\myindex{vizinhança}}).
\end{defn}

\begin{myfigureht}
\myincludepdft{msopenset}{%
Uma região do plano, marcada em cinza-claro,
com fronteira pontilhada e identificada por V.
Dentro da região há um círculo pontilhado, centrado
em um ponto identificado por x, de raio delta e identificado
como B de x e delta. O círculo e seu interior,
sombreados em cinza mais escuro, estão inteiramente dentro da
região em cinza-claro.}
\caption{Conjunto aberto em um espaço métrico. Note que $\delta$ depende de $x$.\label{fig:msopenset}}
\end{myfigureht}

Intuitivamente, um conjunto aberto $V$ é um conjunto que não inclui sua
\myquote{fronteira}.
Em qualquer ponto de $V$, podemos nos \myquote{mover} um pouco e
continuar em $V$. Da mesma forma, um conjunto $E$ é fechado se todos os pontos
que não pertencem a $E$ estão a alguma distância de $E$.
As bolas abertas e fechadas são exemplos de conjuntos abertos e fechados
(isso ainda precisa ser demonstrado).
Mas nem todo conjunto é aberto ou fechado. Em geral, a maioria dos
subconjuntos não é nenhum dos dois.

\begin{example}
O conjunto $(0,\infty) \subset \R$ é aberto: dado qualquer $x \in (0,\infty)$,
seja $\delta \coloneqq x$. Então $B(x,\delta) = (0,2x) \subset (0,\infty)$.

O conjunto $[0,\infty) \subset \R$ é fechado: dado $x \in (-\infty,0) =
[0,\infty)^c$,
seja $\delta \coloneqq -x$. Então $B(x,\delta) = (-2x,0) \subset
(-\infty,0) = [0,\infty)^c$.

O conjunto $[0,1) \subset \R$ não é aberto nem fechado. Primeiro, toda bola
em $\R$ centrada em $0$, $B(0,\delta) = (-\delta,\delta)$, contém números
negativos e, portanto, não está contida em $[0,1)$. Logo, $[0,1)$ não é
aberto. Segundo, toda bola em $\R$ centrada em $1$,
$B(1,\delta) = (1-\delta,1+\delta)$, contém números estritamente menores
que 1 e maiores que 0 (por exemplo, $1-\nicefrac{\delta}{2}$, desde que
$\delta < 2$).
Assim, $[0,1)^c = \R \setminus [0,1)$ não é aberto, e $[0,1)$ não é fechado.

Se $(X,d)$ é um espaço métrico qualquer e $x \in X$, então $\{ x \}$ é
fechado (exercício).
Por outro lado, dependendo de $X$, $\{ x \}$ pode ou não ser aberto.
O conjunto $\{ 0 \} \subset \R$ não é aberto, pois $B(0,\delta)$ contém
números diferentes de zero para todo $\delta > 0$.
Se $X=\{ x \}$, então $\{ x \}$ é aberto.
\end{example}

\begin{prop} \label{prop:topology:open}
Seja $(X,d)$ um espaço métrico.
\begin{enumerate}[(i)]
\item \label{topology:openi} $\emptyset$ e $X$ são abertos.
\item \label{topology:openii} Se $V_1, V_2, \ldots, V_k$ são subconjuntos abertos de $X$, então
\begin{equation*}
\bigcap_{j=1}^k V_j
\end{equation*}
também é aberto. Isto é, uma interseção finita de conjuntos abertos é aberta.
\item \label{topology:openiii} Se $\{ V_\lambda \}_{\lambda \in I}$ é
uma coleção arbitrária de subconjuntos abertos de $X$, então
\begin{equation*}
\bigcup_{\lambda \in I} V_\lambda
\end{equation*}
também é aberto. Isto é, uma união de conjuntos abertos é aberta.
\end{enumerate}
\end{prop}

O conjunto de índices $I$ em \ref{topology:openiii} pode ser arbitrariamente
grande.
Por $\bigcup_{\lambda \in I} V_\lambda$, queremos dizer simplesmente o
conjunto de todos os $x$ tais que $x \in V_\lambda$ para pelo menos um
$\lambda \in I$.

\begin{proof}
Os conjuntos $\emptyset$ e $X$ são obviamente abertos em $X$.

Demonstremos \ref{topology:openii}.
Se $x \in \bigcap_{j=1}^k V_j$, então $x \in V_j$ para todo $j$.
Como todos os $V_j$ são abertos, para cada $j$ existe $\delta_j > 0$ tal que
$B(x,\delta_j) \subset V_j$. Tome $\delta \coloneqq \min \{
\delta_1,\delta_2,\ldots,\delta_k \}$ e observe que $\delta > 0$. Temos
$B(x,\delta) \subset B(x,\delta_j) \subset V_j$ para todo $j$ e, portanto,
$B(x,\delta) \subset \bigcap_{j=1}^k V_j$. Consequentemente, a interseção é
aberta.

Demonstremos \ref{topology:openiii}.
Se $x \in \bigcup_{\lambda \in I} V_\lambda$, então $x \in V_\lambda$ para
algum $\lambda \in I$.
Como $V_\lambda$ é aberto, existe $\delta > 0$ tal que
$B(x,\delta) \subset V_\lambda$. Mas, então,
$B(x,\delta) \subset \bigcup_{\lambda \in I} V_\lambda$,
e, portanto, a união é aberta.
\end{proof}

\begin{example}
Observe a diferença entre os
itens
\ref{topology:openii} e \ref{topology:openiii}.
O item \ref{topology:openii} não vale para uma interseção arbitrária.
Por exemplo,
$\bigcap_{n=1}^\infty (\nicefrac{-1}{n},\nicefrac{1}{n}) = \{ 0 \}$,
que não é aberto.
\end{example}


A demonstração da proposição análoga a seguir, para conjuntos fechados,
fica como exercício.

\begin{prop} \label{prop:topology:closed}
%FIXMEevillayouthack
\pagebreak[2]
Seja $(X,d)$ um espaço métrico.
\begin{enumerate}[(i)]
\item \label{topology:closedi} $\emptyset$ e $X$ são fechados.
\item \label{topology:closedii} Se $\{ E_\lambda \}_{\lambda \in I}$ é
uma coleção arbitrária de subconjuntos fechados de $X$, então
\begin{equation*}
\bigcap_{\lambda \in I} E_\lambda
\end{equation*}
também é fechado. Isto é, uma interseção de conjuntos fechados é fechada.
\item \label{topology:closediii} Se $E_1, E_2, \ldots, E_k$ são
subconjuntos fechados de $X$, então
\begin{equation*}
\bigcup_{j=1}^k E_j
\end{equation*}
também é fechado. Isto é, uma união finita de conjuntos fechados é fechada.
\end{enumerate}
\end{prop}

Apesar dos nomes,
ainda não demonstramos que a bola aberta é aberta e que a bola fechada é
fechada. Demonstremos esses fatos agora para justificar a terminologia.

\begin{prop} \label{prop:topology:ballsopenclosed}
Sejam $(X,d)$ um espaço métrico, $x \in X$ e $\delta > 0$. Então,
$B(x,\delta)$ é aberto e $C(x,\delta)$ é fechado.
\end{prop}

\begin{proof}
Seja $y \in B(x,\delta)$. Seja $\alpha \coloneqq \delta-d(x,y)$. Como
$\alpha > 0$, consideremos $z \in B(y,\alpha)$. Então,
\begin{equation*}
d(x,z) \leq d(x,y) + d(y,z) < d(x,y) + \alpha = d(x,y) + \delta-d(x,y) =
\delta .
\end{equation*}
Portanto, $z \in B(x,\delta)$ para todo $z \in B(y,\alpha)$. Assim,
$B(y,\alpha) \subset B(x,\delta)$ e, portanto,
$B(x,\delta)$ é aberto. Veja a \figureref{fig:ballisopen}.

\begin{myfigureht}
\myincludepdft{ballisopen}{%
Um disco sombreado, com fronteira pontilhada e raio delta, está centrado em x
e identificado como B de x e delta. Outro disco menor está inteiramente
dentro do disco maior, toca sua fronteira por dentro, está centrado em y e tem
raio alpha. Um terceiro ponto z está dentro do disco menor. Os pontos x, y e z
formam um triângulo tracejado.}
\caption{Demonstração de que $B(x,\delta)$ é aberto: $B(y,\alpha) \subset
B(x,\delta)$, com a desigualdade triangular ilustrada.\label{fig:ballisopen}}
\end{myfigureht}

A demonstração de que $C(x,\delta)$ é fechado fica como exercício.
\end{proof}

Novamente, tenha cuidado com o espaço métrico em que estamos.
O conjunto $[0,\nicefrac{1}{2})$ é uma bola aberta em $[0,1]$ e, portanto,
$[0,\nicefrac{1}{2})$ é aberto em $[0,1]$. Por outro lado,
$[0,\nicefrac{1}{2})$ não é aberto nem fechado em $\R$.

\begin{prop} \label{prop:topology:intervals:openclosed}
Sejam $a,b$ números reais, com $a < b$. Então $(a,b)$, $(a,\infty)$ e
$(-\infty,b)$ são abertos em $\R$.
Além disso, $[a,b]$, $[a,\infty)$ e $(-\infty,b]$ são fechados em $\R$.
\end{prop}

A demonstração fica como exercício. Tenha em mente que há muitos outros
conjuntos abertos e fechados nos números reais. %$\R$.

\begin{prop} \label{prop:topology:subspaceopen}
Suponha que $(X,d)$ seja um espaço métrico e $Y \subset X$. Então,
$U \subset Y$ é aberto em $Y$ (na topologia de subespaço) se, e somente se,
existe um conjunto aberto $V \subset X$ (isto é, aberto em $X$) tal que
$V \cap Y = U$.
\end{prop}

Por exemplo, sejam $X \coloneqq \R$, $Y \coloneqq [0,1]$ e
$U \coloneqq [0,\nicefrac{1}{2})$.
Vimos que $U$ é aberto em $Y$.
Podemos tomar $V \coloneqq (\nicefrac{-1}{2},\nicefrac{1}{2})$.

\begin{proof}
Suponha que $V \subset X$ seja aberto e $V \cap Y = U$.
Seja $x \in U$.
Como $V$ é aberto e $x \in V$, existe $\delta > 0$ tal que
$B_X(x,\delta) \subset V$.
Então,
\begin{equation*}
B_Y(x,\delta) = B_X(x,\delta) \cap Y \subset V \cap Y = U .
\end{equation*}
Aí, $U$ é aberto em $Y$.

A demonstração da recíproca, isto é, de que, se $U \subset Y$ é aberto na
topologia de subespaço, então existe um conjunto $V$, fica como
\exerciseref{exercise:mssubspace}.
\end{proof}

Uma dica para concluir a demonstração (o exercício) é pensar em um conjunto
aberto como uma união de bolas abertas. Se $U$ é aberto, então, para cada
$x \in U$, existe $\delta_x > 0$ (dependente de $x$) tal que
$B(x,\delta_x) \subset U$. Então $U = \bigcup_{x\in U} B(x,\delta_x)$.

No caso de um subconjunto aberto de um conjunto aberto ou de um subconjunto
fechado de um conjunto fechado, a situação é mais simples.

\begin{prop} \label{prop:topology:subspacesame}
Suponha que $(X,d)$ seja um espaço métrico, que $V \subset X$ seja aberto
e que $E \subset X$ seja fechado.
\begin{enumerate}[(i)]
\item \label{prop:topology:subspacesame:i}
$U \subset V$ é aberto na topologia de subespaço se, e somente se, $U$ é
aberto em $X$.
\item \label{prop:topology:subspacesame:ii}
$F \subset E$ é fechado na topologia de subespaço se, e somente se, $F$ é
fechado em $X$.
\end{enumerate}
\end{prop}

\begin{proof}
Demonstramos
\ref{prop:topology:subspacesame:i}
e deixamos
\ref{prop:topology:subspacesame:ii} como exercício.

Se $U \subset V$ é aberto na topologia de subespaço, pela
\propref{prop:topology:subspaceopen} existe um conjunto $W \subset X$ aberto
em $X$ tal que $U = W \cap V$. A interseção de dois conjuntos abertos é
aberta; portanto, $U$ é aberto em $X$.

Agora suponha que $U$ seja aberto em $X$. Então $U = U \cap V$. Logo,
$U$ é aberto em $V$, novamente pela \propref{prop:topology:subspaceopen}.
\end{proof}

\subsection{Conjuntos conexos}

Vamos generalizar a ideia de intervalo para espaços métricos em geral. Uma
das principais características de um intervalo em $\R$ é ser conexo: podemos
nos mover continuamente de um ponto a outro sem dar saltos.
Por exemplo, em $\R$ geralmente estudamos funções em intervalos e, em espaços
métricos mais gerais, estudamos funções em conjuntos conexos.

\begin{defn}
Um espaço métrico $(X,d)$ não vazio%
\footnote{Alguns autores não excluem o conjunto vazio da definição; nesse
caso, o conjunto vazio seria conexo.
Nós o excluímos essencialmente pela mesma razão pela qual 1 não é primo nem
composto: nossos conjuntos conexos têm exatamente dois subconjuntos clopen,
enquanto os conjuntos desconexos têm mais de dois. O conjunto vazio tem
exatamente um.}
é \emph{\myindex{conexo}} se os únicos subconjuntos de $X$ que são
simultaneamente abertos e fechados (às vezes chamados de subconjuntos
\emph{\myindex{clopen}}) são $\emptyset$ e o próprio $X$.
Se $(X,d)$ não vazio não é conexo, dizemos que ele é
\emph{\myindex{desconexo}}.

Quando aplicamos o termo \emph{conexo} a um subconjunto não vazio $A \subset X$,
queremos dizer que $A$, com a topologia de subespaço, é conexo.
\end{defn}

Em outras palavras, um conjunto $X$ não vazio é conexo se, sempre que
escrevemos $X = X_1 \cup X_2$, com $X_1 \cap X_2 = \emptyset$ e $X_1$ e
$X_2$ abertos, temos $X_1 = \emptyset$ ou $X_2 = \emptyset$.
Assim, para mostrar que $X$ é desconexo, precisamos encontrar conjuntos
abertos, disjuntos e não vazios $X_1$ e $X_2$ cuja união seja $X$.
Para subconjuntos, enunciamos essa ideia na proposição a seguir.
A proposição está ilustrada na \figureref{fig:disconsubset}.

\begin{prop}
Seja $(X,d)$ um espaço métrico.
Um conjunto $S \subset X$ não vazio é desconexo se, e somente se,
existem conjuntos abertos $U_1$ e $U_2$ em $X$ tais que
$U_1 \cap U_2 \cap S = \emptyset$,
$U_1 \cap S \neq \emptyset$,
$U_2 \cap S \neq \emptyset$, e
\begin{equation*}
S = 
\bigl( U_1 \cap S \bigr)
\cup
\bigl( U_2 \cap S \bigr) .
\end{equation*}
\end{prop}

\begin{myfigureht}
\myincludepdft{disconsubset}{%
Duas regiões sombreadas, cada uma com fronteira contínua,
estão identificadas por S. Uma linha pontilhada contorna cada
região.
O interior da primeira está identificado por U de índice 1,
e o da segunda, por U de índice 2.
Os interiores das duas linhas pontilhadas se intersectam, mas
cada região sombreada está dentro de apenas uma das linhas
pontilhadas.}
\caption{Subconjunto desconexo. Note que $U_1 \cap U_2$ não precisa
ser vazio, mas $U_1 \cap U_2 \cap S = \emptyset$.\label{fig:disconsubset}}
\end{myfigureht}


\begin{proof}
Primeiro, suponha que $S$ seja desconexo: existem conjuntos $S_1$ e $S_2$
não vazios, disjuntos e abertos em $S$, tais que $S = S_1 \cup S_2$.
Pela \propref{prop:topology:subspaceopen}, existem $U_1$ e $U_2$ abertos em
$X$ tais que $U_1 \cap S = S_1$ e $U_2 \cap S = S_2$.

Para a recíproca, comecemos pelos conjuntos $U_1$ e $U_2$.
Então $U_1 \cap S$ e $U_2 \cap S$ são abertos em $S$ pela
\propref{prop:topology:subspaceopen}.
Pela discussão anterior à proposição, $S$ é desconexo.
\end{proof}

\begin{example}
Suponha que $S \subset \R$ e que existam $x,y,z$ tais que $x < z < y$, com
$x,y \in S$ e $z \notin S$. Afirmação: \emph{$S$ é desconexo}. Demonstração:
observe que
\begin{equation*}
\bigl( (-\infty,z) \cap S \bigr)
\cup
\bigl( (z,\infty) \cap S \bigr)
= S .
\end{equation*}
\end{example}

\begin{prop}
Um conjunto $S \subset \R$ não vazio é conexo se, e somente se, $S$ é um
intervalo ou um conjunto unitário.
\end{prop}

\begin{proof}
Suponha que $S$ seja conexo. Se $S$ for um conjunto unitário, terminamos.
Suponha, então, que $x < y$ e $x,y \in S$. Se $z \in \R$ é tal que
$x < z < y$, então $(-\infty,z) \cap S$ e $(z,\infty) \cap S$ são não vazios
e disjuntos. Como $S$ é conexo, a união desses conjuntos não pode ser $S$;
portanto, $z \in S$.
Pela \propref{prop:intervaldef}, $S$ é um intervalo.

Se $S$ é um conjunto unitário, então é conexo.
Suponhamos, portanto, que $S$ seja um intervalo.
Consideremos subconjuntos abertos $U_1$ e $U_2$ de $\R$ tais que
$U_1 \cap S$ e $U_2 \cap S$ sejam não vazios e
$S = 
\bigl( U_1 \cap S \bigr)
\cup
\bigl( U_2 \cap S \bigr)$. Mostraremos que $U_1 \cap S$
e $U_2 \cap S$ têm um ponto em comum e, portanto, não são disjuntos,
o que prova que $S$ é conexo.
Suponha que $x \in U_1 \cap S$ e $y \in U_2 \cap S$.
Sem perda de generalidade, suponha que $x < y$.
Como $S$ é um intervalo, $[x,y] \subset S$.
Observe que $U_2 \cap [x,y] \neq \emptyset$ e defina
$z \coloneqq \inf (U_2 \cap [x,y])$.
Queremos mostrar que $z \in U_1$.
Se $z = x$, então $z \in U_1$.
Se $z > x$, para todo $\epsilon > 0$, a bola
$B(z,\epsilon) = (z-\epsilon,z+\epsilon)$ contém pontos de $[x,y]$ que não
pertencem a $U_2$, pois $z$ é o ínfimo de $U_2 \cap [x,y]$.
Logo, $z \notin U_2$, já que $U_2$ é aberto.
Portanto, $z \in U_1$, pois todo ponto de $[x,y]$ pertence a $U_1$ ou a
$U_2$.
Como $U_1$ é aberto, $B(z,\delta) \subset U_1$ para algum $\delta > 0$
suficientemente pequeno.
Como $z$ é o ínfimo do conjunto não vazio $U_2 \cap [x,y]$, existe
$w \in U_2 \cap [x,y]$ tal que $w \in [z,z+\delta) \subset B(z,\delta)
\subset U_1$.
Assim, $w \in U_1 \cap U_2 \cap [x,y]$.
Logo, $U_1 \cap S$ e $U_2 \cap S$ não são disjuntos, e $S$ é conexo.
Veja a \figureref{fig:intervalcon}.
\end{proof}

\begin{myfigureht}
\myincludepdft{intervalcon}{%
Diagrama da reta real com um ponto x indicado à esquerda e um ponto y à
direita. Entre x e y estão os pontos z e w. Um intervalo aberto que começa
em algum ponto à esquerda e se estende além de w (contendo x, z e w) está
identificado como U de índice 1.
Um intervalo aberto que parte de z e segue até além do diagrama, à direita,
está identificado como U de índice 2.
Um pequeno intervalo aberto entre z menos delta e z mais delta, contendo w,
está indicado.}
\caption{Demonstração de que um intervalo é conexo.\label{fig:intervalcon}}
\end{myfigureht}

\begin{example}
Muitas vezes, uma bola $B(x,\delta)$ é conexa, mas isso não é necessariamente
verdadeiro em todo espaço métrico.
No exemplo mais simples, consideremos um espaço com dois pontos
$\{ a, b\}$ e a métrica discreta.
Então $B(a,2) = \{ a , b \}$ não é conexa, pois $B(a,1) = \{ a \}$ e
$B(b,1) = \{ b \}$ são abertos e disjuntos.
\end{example}

\subsection{Fecho e fronteira}

Às vezes, queremos acrescentar a um conjunto todos os pontos dos quais
podemos nos aproximar a partir dele. Esse conceito é chamado de fecho.

\begin{defn}
Seja $(X,d)$ um espaço métrico e $A \subset X$.
O \emph{\myindex{fecho}} de $A$ é o conjunto
\glsadd{not:closure}
\begin{equation*}
\widebar{A} \coloneqq \bigcap \{ E \subset X : E \text{ é fechado e } A \subset
E \} .
\end{equation*}
\end{defn}

Isto é,
$\widebar{A}$ é a interseção de todos os conjuntos fechados que contêm $A$.

\begin{prop}
Seja $(X,d)$ um espaço métrico e $A \subset X$. O fecho $\widebar{A}$
é fechado, e $A \subset \widebar{A}$.
Além disso, se $A$ é fechado, então $\widebar{A} = A$.
\end{prop}

\begin{proof}
O fecho é uma interseção de conjuntos fechados; portanto, $\widebar{A}$ é
fechado.
Há pelo menos um conjunto fechado que contém $A$, a saber, o próprio $X$,
de modo que $A \subset \widebar{A}$.
Se $A$ é fechado, então $A$ é um conjunto fechado que contém $A$.
Logo, $\widebar{A} \subset A$ e, portanto, $A = \widebar{A}$.
\end{proof}

\begin{example}
O fecho de $(0,1)$ em $\R$ é $[0,1]$. De fato, se
$E$ é fechado e contém $(0,1)$, então $E$ contém $0$ e $1$ (por quê?).
Assim, $[0,1] \subset E$. Mas $[0,1]$ também é fechado.
Logo, o fecho $\overline{(0,1)} = [0,1]$.
\end{example}

\begin{example}
Tenha cuidado para observar em qual espaço métrico ambiente você está
trabalhando.
Se $X = (0,\infty)$, então o fecho de $(0,1)$ em $(0,\infty)$ é $(0,1]$.
De fato, assim como acima, $(0,1]$ é fechado em $(0,\infty)$ (por quê?).
Todo conjunto fechado $E$ que contém $(0,1)$ também contém 1 (por quê?).
Portanto, $(0,1] \subset E$ e, assim,
$\overline{(0,1)} = (0,1]$ quando trabalhamos em $(0,\infty)$.
\end{example}

Justifiquemos a afirmação de que o fecho contém todos os pontos dos quais
podemos nos \myquote{aproximar} a partir do conjunto.

\begin{prop} \label{prop:msclosureappr}
Seja $(X,d)$ um espaço métrico e $A \subset X$. Então $x \in \widebar{A}$
se, e somente se, para todo $\delta > 0$, $B(x,\delta) \cap A \neq \emptyset$.
\end{prop}

\begin{proof}
Demonstremos as duas contrapositivas.
Mostremos que $x \notin \widebar{A}$ se, e somente se, existe
$\delta > 0$ tal que $B(x,\delta) \cap A = \emptyset$.

Primeiro, suponha que $x \notin \widebar{A}$. Sabemos que $\widebar{A}$ é
fechado. Portanto, existe $\delta > 0$ tal que
$B(x,\delta) \subset \widebar{A}^c$. Como $A \subset \widebar{A}$,
temos $B(x,\delta) \subset \widebar{A}^c \subset A^c$ e, consequentemente,
$B(x,\delta) \cap A = \emptyset$.

Por outro lado, suponha que exista $\delta > 0$ tal que
$B(x,\delta) \cap A = \emptyset$.
Em outras palavras,
$A \subset {B(x,\delta)}^c$.
Como ${B(x,\delta)}^c$ é fechado, como $x \not \in {B(x,\delta)}^c$
e como $\widebar{A}$ é a interseção dos conjuntos fechados que contêm $A$,
temos $x \notin \widebar{A}$.
\end{proof}

Também podemos falar do interior de um conjunto (os pontos dos quais não
podemos nos aproximar a partir do complementar) e da fronteira de um
conjunto (os pontos dos quais podemos nos aproximar tanto a partir do
conjunto quanto do seu complementar).

\begin{defn}
Seja $(X,d)$ um espaço métrico e $A \subset X$.
O \emph{\myindex{interior}} de $A$ é o conjunto
\glsadd{not:interior}
\begin{equation*}
A^\circ \coloneqq \{ x \in A : \text{existe } \delta > 0
\text{ tal que } B(x,\delta) \subset A \} .
\end{equation*}
A \emph{\myindex{fronteira}} de $A$ é o conjunto
\glsadd{not:boundary}
\begin{equation*}
\partial A \coloneqq \widebar{A}\setminus A^\circ.
\end{equation*}
\end{defn}

De forma equivalente, o interior é a união dos conjuntos abertos contidos
em $A$; veja \exerciseref{exercise:interiortopologic}. Por definição,
$A^\circ \subset A$; contudo, os pontos da fronteira podem ou não pertencer
a $A$.

\begin{example}
Suponha que $A \coloneqq (0,1]$ e $X \coloneqq \R$. Então
$\widebar{A}=[0,1]$, $A^\circ = (0,1)$ e $\partial A = \{ 0, 1 \}$.
\end{example}

\begin{example}
Considere $X \coloneqq \{ a, b \}$ com a métrica discreta e seja
$A \coloneqq \{ a \}$. Então $\widebar{A} = A^\circ = A$ e
$\partial A = \emptyset$.
\end{example}

\begin{prop}
Seja $(X,d)$ um espaço métrico e $A \subset X$. Então $A^\circ$ é aberto
e $\partial A$ é fechado.
\end{prop}

\begin{proof}
Dado $x \in A^\circ$, existe $\delta > 0$ tal que $B(x,\delta) \subset A$.
Se $z \in B(x,\delta)$, como as bolas abertas são abertas, existe
$\epsilon > 0$ tal que $B(z,\epsilon) \subset B(x,\delta) \subset A$.
Logo, $z \in A^\circ$. Portanto, $B(x,\delta) \subset A^\circ$ e,
assim, $A^\circ$ é aberto.

Como $A^\circ$ é aberto,
$\partial A = \widebar{A} \setminus A^\circ = \widebar{A} \cap
{(A^\circ)}^c$ é fechado.
\end{proof}

A fronteira é o conjunto dos pontos que estão próximos tanto do conjunto
quanto do seu complementar. Veja na \figureref{fig:msboundary} o diagrama
que ilustra a proposição a seguir.

\begin{prop}
Seja $(X,d)$ um espaço métrico e $A \subset X$. Então $x \in \partial A$
se, e somente se, para todo $\delta > 0$,
$B(x,\delta) \cap A$ e
$B(x,\delta) \cap A^c$ são ambos não vazios.
\end{prop}

\begin{myfigureht}
\myincludepdft{msboundary}{%
Uma imagem próxima à fronteira de um conjunto A sombreado.
A fronteira é representada por uma linha contínua, com A abaixo
e o complementar de A acima. Um círculo pontilhado de
raio delta está centrado em um ponto x da fronteira e
está identificado como B de x vírgula delta.}
\caption{A fronteira é o conjunto dos pontos cujas bolas contêm pontos do
conjunto e também do seu complementar.\label{fig:msboundary}}
\end{myfigureht}

\begin{proof}
Suponha que $x \in \partial A = \widebar{A} \setminus A^\circ$ e
seja $\delta > 0$ arbitrário.
Pela \propref{prop:msclosureappr}, $B(x,\delta)$ contém
um ponto de $A$. Se $B(x,\delta)$ não contivesse pontos de $A^c$,
então $x$ pertenceria a $A^\circ$. Portanto, $B(x,\delta)$ também contém
um ponto de $A^c$.

Demonstremos a implicação contrária por contraposição.
Suponha que $x \notin \partial A$, isto é, que $x \notin \widebar{A}$ ou
$x \in A^\circ$.
Se $x \notin \widebar{A}$, então
$B(x,\delta) \subset \widebar{A}^c$
para algum $\delta > 0$, pois $\widebar{A}$ é fechado.
Logo, $B(x,\delta) \cap A$ é vazio, pois $\widebar{A}^c \subset
A^c$.
Se $x \in A^\circ$, então
$B(x,\delta) \subset A$ para algum $\delta > 0$,
de modo que $B(x,\delta) \cap A^c$ é vazio.
\end{proof}

Obtemos o corolário imediato a seguir sobre os fechos de $A$ e $A^c$.
Basta aplicar a \propref{prop:msclosureappr}.

\begin{cor}
Seja $(X,d)$ um espaço métrico e $A \subset X$.
Então $\partial A = \widebar{A} \cap \widebar{A^c}$.
\end{cor}

\subsection{Exercícios}

\begin{exercise}
Demonstre \propref{prop:topology:closed}. Dica:
aplique \propref{prop:topology:open} aos complementares dos conjuntos.
\end{exercise}

\begin{exercise}
Complete a demonstração da \propref{prop:topology:ballsopenclosed}
provando que $C(x,\delta)$ é fechado.
\end{exercise}

\begin{exercise}
Demonstre \propref{prop:topology:intervals:openclosed}.
\end{exercise}

\begin{exercise}
Suponha que $(X,d)$ seja um espaço métrico não vazio com topologia discreta.
Mostre que $X$ é conexo se, e somente se, contém exatamente um elemento.
\end{exercise}

\begin{exercise}
Considere $\Q$ com a métrica usual, $d(x,y) = \sabs{x-y}$, como nosso espaço
métrico.
Demonstre que $\Q$ é
\emph{\myindex{totalmente separado}}, isto é, mostre que, para quaisquer
$x, y \in \Q$ com $x \neq y$, existem dois conjuntos abertos $U$ e $V$ tais
que $x \in U$, $y \in V$, $U \cap V = \emptyset$ e $U \cup V = \Q$.
\end{exercise}

\begin{exercise}
Mostre que, em um espaço métrico,
todo conjunto aberto pode ser escrito como união de conjuntos fechados.
\end{exercise}

\begin{samepage}
\begin{exercise}
Demonstre que, em um espaço métrico,
\begin{enumerate}[a)]
\item
$E$ é fechado se, e somente se, $\partial E \subset E$.
\item
$U$ é aberto se, e somente se, $\partial U \cap U = \emptyset$.
\end{enumerate}
\end{exercise}
\end{samepage}

\begin{exercise}
Demonstre que, em um espaço métrico,
\begin{enumerate}[a)]
\item
$A$ é aberto se, e somente se, $A^\circ = A$.
\item
$U \subset A^\circ$
para todo conjunto aberto $U$ tal que $U \subset A$.
\end{enumerate}
\end{exercise}

\begin{exercise}
Seja $X$ um conjunto, e sejam $d$ e $d'$ duas métricas em $X$.
Suponha que existam $\alpha > 0$ e $\beta > 0$
tais que $\alpha d(x,y) \leq d'(x,y) \leq \beta d(x,y)$ para quaisquer
$x,y \in X$.
Mostre que $U$ é aberto em $(X,d)$ se, e somente se, $U$ é aberto em $(X,d')$.
Isto é, as topologias de $(X,d)$ e $(X,d')$ são iguais.
\end{exercise}

\begin{exercise}
Suponha que $\{ S_i \}$, $i \in \N$,
seja uma coleção de subconjuntos conexos de um espaço métrico $(X,d)$
e que exista $x \in X$ tal que $x \in S_i$ para todo $i \in \N$.
Mostre que $\bigcup_{i=1}^\infty S_i$ é conexo.
\end{exercise}

\begin{exercise}
\pagebreak[2]
Seja $A$ um conjunto conexo em um espaço métrico.
\begin{enumerate}[a)]
\item
O fecho $\widebar{A}$ é conexo? Demonstre ou encontre um contraexemplo.
\item
O interior $A^\circ$ é conexo? Demonstre ou encontre um contraexemplo.
\end{enumerate}
Dica: pense em conjuntos de $\R^2$.
\end{exercise}

\begin{exercise} \label{exercise:mssubspace}
Complete a demonstração da \propref{prop:topology:subspaceopen}.
Suponha que $(X,d)$ seja um espaço métrico e $Y \subset X$. Mostre que,
com a métrica de subespaço em $Y$, se um conjunto $U \subset Y$
é aberto em $Y$, então existe um conjunto aberto $V \subset X$ tal que
$U = V \cap Y$.
\end{exercise}

\begin{samepage}
\begin{exercise}
Seja $(X,d)$ um espaço métrico.
\begin{enumerate}[a)]
\item
Para todo $x \in X$ e $\delta > 0$, mostre que
$\overline{B(x,\delta)} \subset C(x,\delta)$.
\item
É sempre verdade que
$\overline{B(x,\delta)} = C(x,\delta)$? Demonstre ou encontre um
contraexemplo.
\end{enumerate}
\end{exercise}
\end{samepage}

\begin{exercise} \label{exercise:interiortopologic}
Seja $(X,d)$ um espaço métrico e $A \subset X$. Mostre que
$A^\circ = \bigcup \{ V : V \text{ é aberto e } V \subset A \}$.
\end{exercise}

\begin{exercise}
Complete a demonstração da \propref{prop:topology:subspacesame}.
\end{exercise}

\begin{exercise}
Seja $(X,d)$ um espaço métrico. Mostre que existe uma métrica limitada
$d'$ tal que $(X,d')$ tenha os mesmos conjuntos abertos, isto é, que sua
topologia seja a mesma.
\end{exercise}

\begin{samepage}
\begin{exercise}
Seja $(X,d)$ um espaço métrico.
\begin{enumerate}[a)]
\item
Demonstre que, para todo $x \in X$, ou existe $\delta > 0$ tal que
$B(x,\delta) = \{ x \}$, ou $B(x,\delta)$ é infinito para todo $\delta > 0$.
\item
Encontre um exemplo explícito de $(X,d)$, com $X$ infinito, no qual
para todo $\delta > 0$ e todo $x \in X$, a bola $B(x,\delta)$ seja finita.
\item
Encontre um exemplo explícito de $(X,d)$ no qual
para todo $\delta > 0$ e todo $x \in X$, a bola $B(x,\delta)$ seja um
conjunto infinito enumerável.
\item
Demonstre que, se $X$ é não enumerável, então existem $x \in X$ e
$\delta > 0$ tais que $B(x,\delta)$ é não enumerável.
\end{enumerate}
\end{exercise}
\end{samepage}

\begin{exercise}
Para todo $x \in \R^n$ e todo $\delta > 0$, defina o \myquote{retângulo}
$R(x,\delta) \coloneqq
(x_1-\delta,x_1+\delta) \times
(x_2-\delta,x_2+\delta) \times \cdots \times
(x_n-\delta,x_n+\delta)$. Mostre que esses conjuntos geram os mesmos
conjuntos abertos que as bolas na métrica usual. Isto é, mostre que um
conjunto $U \subset \R^n$ é aberto segundo a métrica usual se, e somente se,
para todo ponto $x \in U$, existe $\delta > 0$ tal que
$R(x,\delta) \subset U$.
\end{exercise}

%%%%%%%%%%%%%%%%%%%%%%%%%%%%%%%%%%%%%%%%%%%%%%%%%%%%%%%%%%%%%%%%%%%%%%%%%%%%%%

\sectionnewpage
\section{Sequências e convergência}
\label{sec:metseqs}

%mbxINTROSUBSECTION

\sectionnotes{1 aula}

\subsection{Sequências}

A noção de sequência em um espaço métrico é muito semelhante à de uma
sequência de números reais.
As definições relacionadas são essencialmente as mesmas que as dos números
reais no sentido de \chapterref{seq:chapter}, substituindo-se $\R$ com
a métrica usual $d(x,y)=\sabs{x-y}$
por um espaço métrico arbitrário $(X,d)$.

\begin{defn}
Uma \emph{\myindex{sequência}} em um espaço métrico $(X,d)$ é uma função
$x \colon \N \to X$. Como antes, escrevemos $x_n$ para o $n$-ésimo elemento
da sequência, e usamos a notação a seguir para indicar a sequência inteira:
\glsadd{not:sequence}
\begin{equation*}
\{ x_n \}_{n=1}^\infty .
\end{equation*}

Uma sequência $\{ x_n \}_{n=1}^\infty$ é \emph{limitada}\index{sequência limitada}
se existem um ponto $p \in X$ e $B \in \R$ tais que
\begin{equation*}
d(p,x_n) \leq B \qquad \text{para todo } n \in \N.
\end{equation*}
Em outras palavras, a sequência $\{x_n\}_{n=1}^\infty$ é limitada quando
o conjunto $\{ x_n : n \in \N \}$
é limitado.

Se $\{ n_k \}_{k=1}^\infty$ é uma sequência de números naturais
tal que $n_{k+1} > n_k$ para todo $k$, então
a sequência $\{ x_{n_k} \}_{k=1}^\infty$\glsadd{not:subsequence}
é chamada de
\emph{\myindex{subsequência}} de $\{x_n \}_{n=1}^\infty$.
\end{defn}

De modo semelhante, definimos convergência.
Veja na \figureref{fig:sequence-convergence-metric} uma ilustração
da definição.

\begin{defn}
Dizemos que uma sequência $\{ x_n \}_{n=1}^\infty$ em um espaço métrico
$(X,d)$ \emph{converge} a um ponto
$p \in X$ se, para todo $\epsilon > 0$, existe $M \in \N$ tal que
$d(x_n,p) < \epsilon$ para todo $n \geq M$. Dizemos que o ponto $p$
é um \emph{limite}\index{limite!de uma sequência em um espaço métrico}
de $\{ x_n \}_{n=1}^\infty$.
Demonstraremos que o limite é único e, por isso, escrevemos
\glsadd{not:limseq}
\begin{equation*}
\lim_{n\to \infty} x_n \coloneqq p .
\end{equation*}

Uma sequência que converge é \emph{convergente}\index{convergente!sequência em um espaço
métrico}.
Caso contrário, ela é \emph{divergente}\index{divergente!sequência em um espaço métrico}.
\begin{myfigureht}
\myincludegraphics{sequence-convergence-metric}{%
Um círculo pontilhado de raio epsilon está centrado em p.
Vários pontos do início da sequência são mostrados, desde x de índice 1 até
x de índice 10. Os pontos x de índice 1, x de índice 3, x de índice 4, x de índice 5 e x de índice 6
estão fora do círculo, mas todos os outros estão dentro.}
\caption{Sequência convergente a $p$. São mostrados os primeiros 10 pontos,
e $M=7$ para esse $\epsilon$.\label{fig:sequence-convergence-metric}}
\end{myfigureht}
\end{defn}

Comecemos provando que o limite é único. A demonstração é quase idêntica
(palavra por palavra) à demonstração do mesmo fato para sequências de
números reais, em \propref{prop:limisunique}.
As demonstrações de muitos resultados conhecidos sobre sequências de números
reais podem ser adaptadas para espaços métricos mais gerais.
Nas demonstrações, devemos substituir $\sabs{x-y}$ por $d(x,y)$ e aplicar
corretamente a desigualdade triangular.

\begin{prop} \label{prop:mslimisunique}
Uma sequência convergente em um espaço métrico tem limite único.
\end{prop}

\begin{proof}
%NOTE: should be word for word the same as 2.1.6
Suponha que $\{ x_n \}_{n=1}^\infty$ tenha limites $x$ e $y$.
Escolha $\epsilon > 0$ arbitrário.
Pela definição, encontre $M_1$ tal que, para todo $n \geq M_1$,
$d(x_n,x) < \nicefrac{\epsilon}{2}$. De modo semelhante, encontre $M_2$
tal que, para todo $n \geq M_2$, tenhamos
$d(x_n,y) < \nicefrac{\epsilon}{2}$. Agora escolha $n$ tal que
$n \geq M_1$ e também $n \geq M_2$, e estime:
\begin{equation*}
\begin{split}
d(y,x)
& \leq
d(y,x_n) + d(x_n,x) \\
& <
\frac{\epsilon}{2} + \frac{\epsilon}{2} = \epsilon .
\end{split}
\end{equation*}
Como $d(y,x) < \epsilon$ para todo $\epsilon > 0$, temos $d(x,y) = 0$
e, portanto, $y=x$. Logo, o limite (se existir) é único.
\end{proof}

As demonstrações das proposições a seguir ficam como exercício.

\begin{prop} \label{prop:msconvbound}
Uma sequência convergente em um espaço métrico é limitada.
\end{prop}

\begin{prop} \label{prop:msconvifa}
Uma sequência $\{ x_n \}_{n=1}^\infty$ em um espaço métrico $(X,d)$ converge
a $p \in X$ se, e somente se, existe uma sequência $\{ a_n \}_{n=1}^\infty$
de números reais tal que
\begin{equation*}
d(x_n,p) \leq a_n \quad \text{para todo } n \in \N,
\qquad
\text{e}
\qquad
\lim_{n\to\infty} a_n = 0.
\end{equation*}
\end{prop}

\begin{prop} \label{prop:mssubseq}
Seja $\{ x_n \}_{n=1}^\infty$ uma sequência em um espaço métrico $(X,d)$.
\begin{enumerate}[(i)]
\item Se $\{ x_n \}_{n=1}^\infty$ converge a $p \in X$, então toda
subsequência $\{ x_{n_k} \}_{k=1}^\infty$
converge a $p$.
\item Se, para algum $K \in \N$, a cauda de ordem $K$,
$\{ x_n \}_{n=K+1}^\infty$
converge a $p \in X$, então
 $\{ x_n \}_{n=1}^\infty$ converge a $p$.
\end{enumerate}
\end{prop}

\begin{example}
Considere $C\bigl([a,b],\R\bigr)$ como o conjunto das funções contínuas,
com a norma uniforme como métrica. Vimos que assim obtemos um espaço métrico.
Pela definição de convergência, notamos que ela é idêntica à convergência
uniforme. Isto é, $\{ f_n \}_{n=1}^\infty$ converge uniformemente se, e
somente se, converge no sentido de espaço métrico.
\end{example}

\begin{remark}
Talvez seja surpreendente que, no conjunto das funções
$f \colon [a,b] \to \R$ (contínuas ou não), não exista uma métrica que
induza a convergência pontual, pois o domínio $[a,b]$ é não enumerável.
A demonstração desse fato está além do escopo deste livro.
\end{remark}

\subsection{Convergência no espaço euclidiano}

No espaço euclidiano $\R^n$, uma sequência converge se, e somente se,
cada coordenada converge:

\begin{prop} \label{prop:msconveuc}
Seja $\{ x_m \}_{m=1}^\infty$ uma sequência em $\R^n$,
em que $x_m = \bigl(x_{m,1},x_{m,2},\ldots,x_{m,n}\bigr) \in \R^n$.
Então $\{ x_m \}_{m=1}^\infty$ converge se, e somente se,
$\{ x_{m,k} \}_{m=1}^\infty$ converge para todo $k=1,2,\ldots,n$; nesse caso,
\begin{equation*}
\lim_{m\to\infty}
x_m =
\Bigl(
\lim_{m\to\infty} x_{m,1},
\lim_{m\to\infty} x_{m,2},
\ldots,
\lim_{m\to\infty} x_{m,n}
\Bigr) .
\end{equation*}
\end{prop}

\begin{proof}
Suponha que
$\{ x_m \}_{m=1}^\infty$ convirja para
$y = (y_1,y_2,\ldots,y_n) \in \R^n$.
Dado $\epsilon > 0$, existe $M$ tal que, para todo
$m \geq M$, temos
\begin{equation*}
d(y,x_m) < \epsilon.
\end{equation*}
Fixe um $k=1,2,\ldots,n$. Para todo $m \geq M$,
\begin{equation*}
\bigl\lvert y_k - x_{m,k} \bigr\rvert
=
\sqrt{{\bigl(y_k - x_{m,k} \bigr)}^2}
\leq
\sqrt{\sum_{\ell=1}^n {\bigl(y_\ell-x_{m,\ell}\bigr)}^2}
= d(y,x_m) < \epsilon .
\end{equation*}
Logo, a sequência $\{ x_{m,k} \}_{m=1}^\infty$ converge para $y_k$.

Para a implicação contrária, suponha que 
$\{ x_{m,k} \}_{m=1}^\infty$ convirja para $y_k$ para todo $k=1,2,\ldots,n$.
Dado $\epsilon > 0$, escolha $M$ tal que, se $m \geq M$, então 
$\bigl\lvert y_k-x_{m,k} \bigr\rvert < \nicefrac{\epsilon}{\sqrt{n}}$ para todo
$k=1,2,\ldots,n$. Então
\begin{equation*}
d(y,x_m)
=
\sqrt{\sum_{k=1}^n {\bigl(y_k-x_{m,k}\bigr)}^2}
<
\sqrt{\sum_{k=1}^n {\left(\frac{\epsilon}{\sqrt{n}}\right)}^2}
=
\sqrt{\sum_{k=1}^n \frac{{\epsilon^2}}{n}}
= \epsilon .
\end{equation*}
Isto é, a sequência $\{ x_m \}_{m=1}^\infty$ converge para
$y = (y_1,y_2,\ldots,y_n) \in \R^n$.
\end{proof}

\begin{example}
Como dissemos, o conjunto $\C$ dos números complexos $z = x+iy$ é considerado
como o espaço métrico $\R^2$. A proposição afirma que a sequência
$\{ z_n \}_{n=1}^\infty = \{ x_n + iy_n \}_{n=1}^\infty$ converge
a $z = x+iy$
se, e somente se, $\{ x_n \}_{n=1}^\infty$ converge para $x$ e
$\{ y_n \}_{n=1}^\infty$ converge para $y$.
\end{example}

\subsection{Convergência e topologia}

A topologia — a coleção dos conjuntos abertos de um espaço — codifica quais
sequências convergem.

\begin{prop} \label{prop:msconvtopo}
Seja $(X,d)$ um espaço métrico e $\{x_n\}_{n=1}^\infty$ uma sequência em $X$.
Então
$\{ x_n \}_{n=1}^\infty$ converge a $p \in X$ se, e somente se, para toda
vizinhança aberta $U$ de $p$, existe $M \in \N$ tal que, para todo $n \geq M$,
temos $x_n \in U$.
\end{prop}

\begin{proof}
Suponha que $\{ x_n \}_{n=1}^\infty$ convirja a $p$. Seja $U$ uma vizinhança
aberta de $p$. Existe $\epsilon > 0$ tal que $B(p,\epsilon) \subset
U$. Como a sequência converge, encontre $M \in \N$ tal que, para todo $n \geq
M$, temos $d(p,x_n) < \epsilon$; em outras palavras,
$x_n \in B(p,\epsilon) \subset U$.

Demonstremos a implicação contrária. Dado $\epsilon > 0$, seja
$U \coloneqq B(p,\epsilon)$ a vizinhança de $p$. Então existe $M \in \N$
tal que, para $n \geq M$, temos $x_n \in U = B(p,\epsilon)$; em outras
palavras, $d(p,x_n) < \epsilon$.
\end{proof}

Um conjunto fechado contém os limites de suas sequências convergentes.

\begin{prop} \label{prop:msclosedlim}
Seja $(X,d)$ um espaço métrico, seja $E \subset X$ um conjunto fechado
e seja $\{ x_n \}_{n=1}^\infty$ uma sequência em $E$ que converge a algum
$p \in X$. Então $p \in E$.
\end{prop}

\begin{proof}
Demonstremos a contrapositiva.
Suponha que $\{ x_n \}_{n=1}^\infty$ seja uma sequência em $X$ que converge
a $p \in E^c$.
Como $E^c$ é aberto, a \propref{prop:msconvtopo} afirma que existe
$M$ tal que, para todo $n \geq M$,
$x_n \in E^c$. Portanto, $\{ x_n \}_{n=1}^\infty$ não é uma sequência em $E$.
\end{proof}

Para obter o fecho de um conjunto $A$, começamos com $A$ e acrescentamos
os pontos que são limites de sequências em $A$.

\begin{prop} \label{prop:msclosureapprseq}
Seja $(X,d)$ um espaço métrico e $A \subset X$.
Então $p \in \widebar{A}$ se, e somente se, existe uma sequência
$\{ x_n \}_{n=1}^\infty$ de elementos de $A$ tal que
$\lim\limits_{n\to\infty} x_n = p$.
\end{prop}

\begin{proof}
Seja $p \in \widebar{A}$. Para todo $n \in \N$,
a \propref{prop:msclosureappr} garante a existência
de um ponto $x_n \in B(p,\nicefrac{1}{n}) \cap A$.
Como $d(p,x_n) < \nicefrac{1}{n}$, temos $\lim_{n\to\infty} x_n = p$.

Para a implicação contrária, veja \exerciseref{exercise:reverseclosedseq}.
\end{proof}

\subsection{Exercícios}

\begin{exercise} \label{exercise:reverseclosedseq}
Complete a demonstração da 
\propref{prop:msclosureapprseq}:
seja $(X,d)$ um espaço métrico e
$A \subset X$. Seja $p \in X$ tal
que existe uma sequência $\{ x_n \}_{n=1}^\infty$ em $A$ que converge a $p$.
Demonstre que $p \in \widebar{A}$.
\end{exercise}

\begin{exercise}
\leavevmode
\begin{enumerate}[a)]
\item
Mostre que $d(x,y) \coloneqq \min \bigl\{ 1, \sabs{x-y} \bigr\}$ define uma
métrica em $\R$.
\item
Mostre que uma sequência converge em $(\R,d)$ se, e somente se, converge
na métrica usual.
\item
Encontre uma sequência limitada em $(\R,d)$ que
não contenha nenhuma subsequência convergente.
\end{enumerate}
\end{exercise}

\begin{exercise}
Demonstre \propref{prop:msconvbound}.
\end{exercise}

\begin{exercise}
Demonstre \propref{prop:msconvifa}.
\end{exercise}

\begin{exercise}
Suponha que $\{x_n\}_{n=1}^\infty$ convirja a $p$. Suponha que
$f \colon \N \to \N$ seja uma função injetiva. Mostre que
$\{ x_{f(n)} \}_{n=1}^\infty$ converge a $p$.
\end{exercise}

\begin{exercise}
Seja $(X,d)$ um espaço métrico em que $d$ é a métrica discreta. Suponha que
$\{ x_n \}_{n=1}^\infty$ seja uma sequência convergente em $X$. Mostre que
existe $K \in \N$ tal que, para todo $n \geq K$, temos $x_n = x_K$.
\end{exercise}

\begin{exercise}
Dizemos que um conjunto $S \subset X$ é \emph{\myindex{denso}} em $X$ se
$X \subset \widebar{S}$; em outras palavras, se, para todo $p \in X$,
existe uma sequência $\{ x_n \}_{n=1}^\infty$ em $S$ que converge a $p$.
Demonstre que $\R^n$ contém um subconjunto enumerável e denso.
\end{exercise}

\begin{exercise}[Desafio]
Suponha que $\{ U_n \}_{n=1}^\infty$ seja uma sequência decrescente de
conjuntos abertos (isto é, $U_{n+1} \subset U_n$ para
todo $n$) em um espaço métrico $(X,d)$, e que
$\bigcap_{n=1}^\infty U_n = \{ p \}$ para algum $p \in X$.
Suponha que $\{ x_n \}_{n=1}^\infty$ seja uma sequência de pontos de $X$
tal que $x_n \in U_n$. A sequência $\{ x_n \}_{n=1}^\infty$ necessariamente
converge a $p$? Demonstre ou construa um contraexemplo.
\end{exercise}

\begin{exercise}
Seja $E \subset X$ fechado e seja
$\{ x_n \}_{n=1}^\infty$ uma sequência em $X$ que converge a $p \in X$.
Suponha que $x_n \in E$ para infinitos $n \in \N$. Mostre que $p \in E$.
\end{exercise}

\begin{exercise} \label{exercise:extendedrealsmetric}
Seja $\R^* = \{ -\infty \} \cup \R \cup \{ \infty \}$ o conjunto dos reais
estendidos.
Defina
$d(x,y) \coloneqq \babs{ \frac{x}{1+\sabs{x}} - \frac{y}{1+\sabs{y}} }$
se $x,y \in \R$; defina
$d(\infty,x) \coloneqq \babs{ 1 - \frac{x}{1+\sabs{x}} }$,
$d(-\infty,x) \coloneqq \babs{ 1 + \frac{x}{1+\sabs{x}} }$
para todo $x \in \R$; e
seja $d(\infty,-\infty) \coloneqq 2$.
\begin{enumerate}[a)]
\item
Mostre que $(\R^*,d)$ é um espaço métrico.
\item
Suponha que $\{ x_n \}_{n=1}^\infty$ seja uma sequência de números reais tal
que, para todo $M \in \R$, existe $N$ tal que
$x_n \geq M$ para todo $n \geq N$. Mostre que
$\lim\limits_{n\to\infty} x_n = \infty$ em
$(\R^*,d)$.
\item
Mostre que uma sequência de números reais converge a um número real
em $(\R^*,d)$ se, e somente se, converge em $\R$ com a métrica usual.
\end{enumerate}
\end{exercise}

\begin{exercise}
Suponha que $\{ V_n \}_{n=1}^\infty$ seja uma sequência de conjuntos abertos
em $(X,d)$
tal que $V_{n+1} \supset V_n$ para todo $n$. Seja $\{ x_n \}_{n=1}^\infty$
uma sequência
tal que $x_n \in V_{n+1} \setminus V_n$ e suponha que
$\{ x_n \}_{n=1}^\infty$ convirja a $p \in X$. Mostre que $p \in \partial V$,
em que $V = \bigcup_{n=1}^\infty V_n$.
\end{exercise}

\begin{exercise}
Demonstre \propref{prop:mssubseq}.
\end{exercise}

\begin{exercise}
Seja $(X,d)$ um espaço métrico e $\{ x_n \}_{n=1}^\infty$ uma sequência em $X$.
Demonstre que $\{ x_n \}_{n=1}^\infty$ converge a $p \in X$
se, e somente se, toda subsequência de $\{ x_n \}_{n=1}^\infty$ possui
uma subsequência que converge a $p$.
\end{exercise}

\begin{exercise}
Considere $\R^n$ e seja $d$ a métrica euclidiana usual.
Sejam
$d'(x,y) \coloneqq \sum_{\ell=1}^n \sabs{x_\ell-y_\ell}$ e
$d''(x,y) \coloneqq \max \bigl\{
\sabs{x_1-y_1},\sabs{x_2-y_2},\cdots,\sabs{x_n-y_n}
\bigr\}$.
\begin{enumerate}[a)]
\item
Use \exerciseref{exercise:mscross} para mostrar que
$(\R^n,d')$ e
$(\R^n,d'')$ são espaços métricos.
\item
Seja $\{ x_k \}_{k=1}^\infty$ uma sequência em $\R^n$ e seja $p \in \R^n$.
Demonstre que as afirmações a seguir são equivalentes:
\begin{enumerate}[(1)]
\item
$\{ x_k \}_{k=1}^\infty$ converge a $p$ em $(\R^n,d)$.
\item
$\{ x_k \}_{k=1}^\infty$ converge a $p$ em $(\R^n,d')$.
\item
$\{ x_k \}_{k=1}^\infty$ converge a $p$ em $(\R^n,d'')$.
\end{enumerate}
\end{enumerate}
\end{exercise}

%%%%%%%%%%%%%%%%%%%%%%%%%%%%%%%%%%%%%%%%%%%%%%%%%%%%%%%%%%%%%%%%%%%%%%%%%%%%%%

\sectionnewpage
\section{Completude e compacidade}
\label{sec:metcompact}

%mbxINTROSUBSECTION

\sectionnotes{2 aulas}

\subsection{Sequências de Cauchy e completude}

Assim como no caso das sequências de números reais, definimos sequências
de Cauchy.

\begin{defn}
Seja $(X,d)$ um espaço métrico.
Uma sequência $\{ x_n \}_{n=1}^\infty$ em $X$ é uma
\emph{\myindex{sequência de Cauchy}} se, para todo $\epsilon > 0$, existe
$M \in \N$ tal que, para todo $n \geq M$ e todo $k \geq M$, temos
\begin{equation*}
d(x_n, x_k) < \epsilon .
\end{equation*}
\end{defn}

Mais uma vez, a definição é apenas uma tradução do conceito dos números
reais para espaços métricos. Uma sequência de números reais é de Cauchy
no sentido de \chapterref{seq:chapter} se, e somente se, é de Cauchy no
sentido acima, desde que equipemos os números reais com a métrica usual
$d(x,y) = \sabs{x-y}$.

\begin{prop}
Uma sequência convergente em um espaço métrico é de Cauchy.
\end{prop}

\begin{proof}
Suponha que $\{ x_n \}_{n=1}^\infty$ convirja a $p$.
Dado $\epsilon > 0$, existe $M$ tal que, para todo $n \geq M$,
temos $d(p,x_n) < \nicefrac{\epsilon}{2}$. Portanto,
para quaisquer $n,k \geq M$, temos
$d(x_n,x_k) \leq d(x_n,p) + d(p,x_k) < \nicefrac{\epsilon}{2} +
\nicefrac{\epsilon}{2} = \epsilon$.
\end{proof}

\begin{defn}
Dizemos que um espaço métrico $(X,d)$ é
\emph{\myindex{completo}} ou \emph{\myindex{completo de Cauchy}}
se toda sequência de Cauchy $\{ x_n \}_{n=1}^\infty$ em $X$
converge a algum $p \in X$.
\end{defn}

\begin{prop}
O espaço $\R^n$ com a métrica usual é um espaço métrico completo.
\end{prop}

A completude de $\R = \R^1$ foi demonstrada em \chapterref{seq:chapter}.
A demonstração da completude de $\R^n$
se reduz ao caso unidimensional.

\begin{proof}
Seja $\{ x_m \}_{m=1}^\infty$ uma sequência de Cauchy
em $\R^n$, em que $x_m = \bigl(x_{m,1},x_{m,2},\ldots,x_{m,n}\bigr) \in \R^n$.
Como a sequência é de Cauchy, dado $\epsilon > 0$, existe $M$ tal que, para
quaisquer $i,j \geq M$,
\begin{equation*}
d(x_i,x_j) < \epsilon.
\end{equation*}

Fixe um $k=1,2,\ldots,n$. Para $i,j \geq M$,
\begin{equation*}
\bigl\lvert x_{i,k} - x_{j,k} \bigr\rvert
=
\sqrt{{\bigl(x_{i,k} - x_{j,k}\bigr)}^2}
\leq
\sqrt{\sum_{\ell=1}^n {\bigl(x_{i,\ell}-x_{j,\ell}\bigr)}^2}
= d(x_i,x_j) < \epsilon .
\end{equation*}
Logo, a sequência $\{ x_{m,k} \}_{m=1}^\infty$ é de Cauchy. Como $\R$ é
completo, a sequência converge; existe $y_k \in \R$ tal que
$y_k = \lim_{m\to\infty} x_{m,k}$.
Escreva $y = (y_1,y_2,\ldots,y_n) \in \R^n$.
Pela \propref{prop:msconveuc}, a sequência $\{ x_m \}_{m=1}^\infty$ converge
a $y \in \R^n$; portanto, $\R^n$ é completo.
\end{proof}

Na linguagem dos espaços métricos,
os resultados sobre continuidade de \sectionref{sec:liminter}
dizem que o espaço métrico
$C\bigl([a,b],\R\bigr)$ de \exampleref{example:msC01} é completo.
A demonstração, que consiste em \myquote{expandir as definições},
fica como exercício em \exerciseref{exercise:CabRcomplete}.

\begin{prop} \label{prop:CabRcomplete}
O espaço das funções contínuas $C\bigl([a,b],\R\bigr)$, com a norma
uniforme como métrica, é um espaço métrico completo.
\end{prop}

Um subconjunto de um espaço métrico completo, como $\R^n$ com a métrica de
subespaço, não precisa ser completo. Por exemplo, $(0,1]$ com a métrica de
subespaço não é completo, pois $\{ \nicefrac{1}{n} \}_{n=1}^\infty$ é uma
sequência de Cauchy em $(0,1]$ sem limite em $(0,1]$.
No entanto, um subespaço fechado de um espaço métrico completo é completo.
Afinal, podemos pensar em um conjunto fechado como aquele que contém todos
os pontos que são limites de sequências nele contidas.
A demonstração fica como exercício em \exerciseref{exercise:closedcomplete}.

\begin{prop} \label{prop:closedcomplete}
Suponha que $(X,d)$ seja um espaço métrico completo e que $E \subset X$
seja fechado.
Então $E$ é um espaço métrico completo com a métrica de subespaço.
\end{prop}

\subsection{Compacidade}

\begin{defn}
Seja $(X,d)$ um espaço métrico e $K \subset X$. 
O conjunto $K$ é \emph{\myindex{compacto}}
se, para toda coleção
de conjuntos abertos $\{ U_{\lambda} \}_{\lambda \in I}$ tal que
\begin{equation*}
K \subset \bigcup_{\lambda \in I} U_\lambda ,
\end{equation*}
existe um subconjunto finito
$\{ \lambda_1, \lambda_2,\ldots,\lambda_m \} \subset I$
tal que
\begin{equation*}
K \subset \bigcup_{j=1}^m U_{\lambda_j} .
\end{equation*}
\end{defn}

Uma coleção de conjuntos abertos $\{ U_{\lambda} \}_{\lambda \in I}$ como
a anterior é chamada de \emph{\myindex{cobertura aberta}} de $K$.
Podemos dizer que $K$ é compacto quando \emph{toda cobertura aberta de $K$
tem uma \myindex{subcobertura} finita}.

\begin{example}
Considere $\R$ como espaço métrico com a métrica usual.

O conjunto $\R$ não é compacto. Demonstração: para $j \in \N$, seja
$U_j \coloneqq (-j,j)$.
Todo $x \in \R$ pertence a algum $U_j$ (pela
\hyperref[thm:arch:i]{propriedade arquimediana}), portanto, temos uma
cobertura aberta.
Suponha que exista uma subcobertura finita
$\R \subset U_{j_1} \cup U_{j_2} \cup \cdots \cup U_{j_m}$,
com $j_1 < j_2 < \cdots < j_m$. Então $\R \subset U_{j_m}$, o que é
contraditório, pois $j_m \in \R$, mas $j_m \notin U_{j_m} =
(-j_m,j_m)$.

O conjunto $(0,1) \subset \R$ também não é compacto. Demonstração: tome os
conjuntos $U_{j} \coloneqq (\nicefrac{1}{j},1-\nicefrac{1}{j})$ para
$j=3,4,5,\ldots$.
Como acima, $(0,1) = \bigcup_{j=3}^\infty U_j$. De modo semelhante, se
existisse uma subcobertura finita, haveria um conjunto $U_j$ tal que
$(0,1) \subset U_j$, o que novamente levaria a uma contradição.

O conjunto $\{ 0 \} \subset \R$ é compacto. Demonstração: dada uma cobertura
aberta $\{ U_{\lambda} \}_{\lambda \in I}$, deve existir $\lambda_0 \in I$
tal que $0 \in U_{\lambda_0}$, pois $\{ U_{\lambda} \}_{\lambda \in I}$ é uma
cobertura.
O único conjunto $U_{\lambda_0}$ forma uma subcobertura finita.

Demonstraremos abaixo que $[0,1]$, e de fato todo intervalo fechado e
limitado $[a,b]$, é compacto.
\end{example}

\begin{prop}
Seja $(X,d)$ um espaço métrico. Se $K \subset X$ é compacto, então
$K$ é fechado e limitado.
\end{prop}

\begin{proof}
Primeiro, provaremos que um conjunto compacto é limitado.
Fixe $p \in X$. Temos a cobertura aberta
\begin{equation*}
K \subset \bigcup_{n=1}^\infty B(p,n) = X .
\end{equation*}
Se $K$ é compacto, existem índices
$n_1 < n_2 < \ldots < n_m$ tais que
\begin{equation*}
K \subset \bigcup_{j=1}^m B(p,n_j) = B(p,n_m) .
\end{equation*}
Como $K$ está contido em uma bola, $K$ é limitado.
Veja o lado esquerdo da \figureref{fig:compactbndclosed}.

Agora mostraremos que um conjunto que não é fechado não é compacto. Suponha
que $\widebar{K} \neq K$, isto é, que exista um ponto $x \in \widebar{K}
\setminus K$.
Se $y \neq x$, então
$y \notin C(x,\nicefrac{1}{n})$
para $n \in \N$
tal que $\nicefrac{1}{n} < d(x,y)$.
Além disso, $x \notin K$, logo
\begin{equation*}
K \subset \bigcup_{n=1}^\infty {C(x,\nicefrac{1}{n})}^c .
\end{equation*}
Uma bola fechada é fechada, portanto, seu complementar ${C(x,\nicefrac{1}{n})}^c$
é aberto e temos uma cobertura aberta.
Se tomarmos qualquer
coleção finita de índices $n_1 < n_2 < \ldots < n_m$, então
\begin{equation*}
\bigcup_{j=1}^m {C(x,\nicefrac{1}{n_j})}^c 
=
{C(x,\nicefrac{1}{n_m})}^c .
\end{equation*}
Como $x$ pertence ao fecho de $K$,
temos
$C(x,\nicefrac{1}{n_m}) \cap K \neq \emptyset$. Portanto, não há
subcobertura finita e $K$ não é compacto.
Veja o lado direito da \figureref{fig:compactbndclosed}.
\end{proof}

\begin{myfigureht}
\myincludepdft{compact_bnd_closed}{%
Dois diagramas. Um conjunto K aparece como região sombreada, delimitada por
uma linha contínua. No diagrama da esquerda, um ponto p está destacado, junto com
três bolas abertas encaixadas (círculos neste caso), de raios 1, 2 e 3,
todas centradas em p e com fronteiras pontilhadas. O conjunto K está
contido na maior delas.
No diagrama da direita, o fundo está sombreado.
Um ponto x na fronteira de K está indicado, e quatro bolas fechadas estão
desenhadas como círculos, com o interior sombreado em tons de cinza cada vez
mais claros, indicando que consideramos os complementares.
A bola externa tem raio 1; seguem-se bolas de raios um meio, um terço e um
quarto. Como x está na fronteira, os complementares das bolas cobrem porções
cada vez maiores de K. Por outro lado, K também tem interseção não vazia com
todas as bolas fechadas.}
\caption{Demonstração de que um conjunto compacto é limitado (à esquerda)
e fechado (à direita).\label{fig:compactbndclosed}}
\end{myfigureht}

Demonstraremos abaixo que, em um espaço euclidiano de dimensão finita,
todo conjunto fechado e limitado é compacto. Portanto, os conjuntos fechados
e limitados de $\R^n$ são exemplos de conjuntos compactos.
Isso não vale em todo espaço métrico: ser fechado e limitado não é
equivalente a ser compacto. Um exemplo simples é o espaço métrico incompleto
$(0,1)$ com a métrica de subespaço herdada de $\R$.
Há muitos espaços métricos completos e muito úteis nos quais ser fechado e
limitado não basta para garantir compacidade: $C\bigl([a,b],\R\bigr)$ é um
espaço métrico completo, mas a bola unitária fechada $C(0,1)$ não é compacta;
veja \exerciseref{exercise:msclbounnotcompt}. Consulte também
\exerciseref{exercise:mstotbound}.

Deixando a limitação de lado por um instante, observe também a diferença
entre ser fechado e ser compacto. Ser fechado depende do espaço métrico
ambiente: o conjunto $(0,1]$ não é fechado em $\R$, mas é fechado no
subespaço $(0,\infty)$.
No entanto, um conjunto $K$ é compacto em um espaço métrico $(X,d)$ se, e
somente se, é compacto com a métrica de subespaço em $K$. Assim, para um
conjunto compacto, não precisamos perguntar em qual espaço métrico ele está.
Por outro lado, todo conjunto é sempre fechado na métrica de subespaço,
considerado como subconjunto de si mesmo.
Veja também \exerciseref{exercise:subspacecompact}.

Uma propriedade útil dos conjuntos compactos em um espaço métrico é que toda
sequência no conjunto tem uma subsequência convergente a um ponto do próprio
conjunto.
Esses conjuntos são chamados de
\emph{\myindex{sequencialmente compacto}}.
Demonstraremos que, em espaços métricos, um conjunto é compacto se, e
somente se, é sequencialmente compacto.
Primeiro, provaremos um lema.

\begin{lemma}[Lema da cobertura de Lebesgue%
\footnote{Nomeado em homenagem ao matemático francês
\href{https://en.wikipedia.org/wiki/Henri_Lebesgue}{Henri L\'eon Lebesgue}
(1875--1941).
O número $\delta$ às vezes é chamado de \myindex{número de Lebesgue} da
cobertura.}]\label{ms:lebesgue}
\index{lema da cobertura de Lebesgue}
Seja $(X,d)$ um espaço métrico e $K \subset X$. Suponha que toda sequência
em $K$ tenha uma subsequência convergente em $K$. Dada
uma cobertura aberta $\{ U_\lambda \}_{\lambda \in I}$ de $K$, existe
$\delta > 0$ tal que, para todo $x \in K$, existe $\lambda \in I$
com $B(x,\delta) \subset U_\lambda$.
\end{lemma}

\begin{proof}
Demonstraremos o lema por contraposição.
Se a conclusão não for verdadeira, então existe
uma cobertura aberta $\{ U_\lambda \}_{\lambda \in I}$ de $K$ com
a seguinte propriedade.
Para todo $n \in \N$, existe $x_n \in K$ tal que
$B(x_n,\nicefrac{1}{n})$ não está contida em nenhum $U_\lambda$.
Tome qualquer $x \in K$. Existe
$\lambda \in I$ tal que $x \in U_\lambda$. Como $U_\lambda$ é aberto,
existe $\epsilon > 0$
tal que $B(x,\epsilon) \subset U_\lambda$. Escolha $M$ tal que
$\nicefrac{1}{M} < \nicefrac{\epsilon}{2}$. Se $y \in
B(x,\nicefrac{\epsilon}{2})$ e $n \geq M$, então
\begin{equation*}
B(y,\nicefrac{1}{n}) \subset
B(y,\nicefrac{1}{M}) \subset
B(y,\nicefrac{\epsilon}{2}) \subset B(x,\epsilon)
\subset U_\lambda ,
\end{equation*}
em que
$B(y,\nicefrac{\epsilon}{2}) \subset B(x,\epsilon)$
decorre da desigualdade triangular.
Veja a \figureref{fig:lebesguedelta}.
Assim, $y \neq x_n$.
Em outras palavras, para todo $n \geq M$, $x_n \notin B(x,\nicefrac{\epsilon}{2})$.
A sequência não pode ter uma subsequência que convirja a $x$. Como $x \in K$
era arbitrário, concluímos a demonstração.
\end{proof}

\begin{myfigureht}
\myincludepdft{lebesguedelta}{%
Um diagrama mostra um conjunto aberto U de índice lambda, sombreado.
Um ponto x nesse conjunto está marcado, e uma bola de raio
epsilon centrada em x está desenhada dentro dele.
Também está desenhada uma bola de raio epsilon dividido por 2 centrada em x,
contida na bola maior.
Um ponto y na bola menor está marcado, e uma bola de raio epsilon dividido por 2
centrada em y está representada. Essa bola centrada em y está inteiramente
dentro da bola de raio epsilon centrada em x, a maior, e portanto também
dentro do conjunto aberto. Por fim, aparece uma pequena bola de raio 1 dividido por n
centrada em y, contida na bola de raio epsilon dividido por 2 centrada em y.}
\caption{Demonstração do lema da cobertura de Lebesgue.
Note que $B(y,\nicefrac{\epsilon}{2}) \subset
B(x,\epsilon)$ pela desigualdade triangular.\label{fig:lebesguedelta}}
\end{myfigureht}

É importante compreender o que o lema afirma. Se $K$ é sequencialmente
compacto, então, para qualquer cobertura dada, existe um mesmo
$\delta > 0$ que serve para todos os pontos. O valor de $\delta$ depende
da cobertura, mas, naturalmente, não depende de $x$.

Por exemplo, seja $K \coloneqq [-10,10]$ e considere a cobertura aberta
formada por $U_n \coloneqq (n,n+2)$, para $n \in \Z$.
Tome $x \in K$. Existe $n \in \Z$ tal que $n \leq x < n+1$.
Se $n \leq x < n+\nicefrac{1}{2}$, então
$B\bigl(x,\nicefrac{1}{2}\bigr) \subset U_{n-1}$.
Se $n+ \nicefrac{1}{2} \leq x < n+1$, então
$B\bigl(x,\nicefrac{1}{2}\bigr) \subset U_{n}$. Portanto,
$\delta = \nicefrac{1}{2}$ serve. Os conjuntos
$U'_n \coloneqq \bigl(\frac{n}{2},\frac{n+2}{2} \bigr)$
também formam uma cobertura aberta,
mas agora o maior valor de $\delta$ que serve é $\nicefrac{1}{4}$.

Por outro lado, $\N \subset \R$ não é sequencialmente compacto.
Fica como exercício encontrar uma cobertura para a qual nenhum
$\delta > 0$ sirva.

\begin{thm} \label{thm:mscompactisseqcpt}
Seja $(X,d)$ um espaço métrico. Então $K \subset X$ é compacto se, e
somente se, toda sequência em $K$ tem uma subsequência que converge a
um ponto de $K$.
\end{thm}

\begin{proof}
Afirmação: \emph{Seja $K \subset X$ um subconjunto de $X$ e
$\{ x_n \}_{n=1}^\infty$ uma sequência em $K$. Suponha que, para cada
$x \in K$, exista uma bola $B(x,\alpha_x)$, com $\alpha_x > 0$, tal que
$x_n \in B(x,\alpha_x)$ para apenas um número finito de $n \in \N$.
Então $K$ não é compacto.}

Demonstração da afirmação:
Observe que
\begin{equation*}
K \subset \bigcup_{x \in K} B(x,\alpha_x) .
\end{equation*}
Qualquer coleção finita dessas bolas contém apenas um número finito de
elementos da sequência $\{ x_n \}_{n=1}^\infty$,
portanto, existe um $x_n \in K$
que não pertence à união. Logo, $K$ não é compacto, e a afirmação está
demonstrada.

Suponha agora que $K$ seja compacto e que $\{ x_n \}_{n=1}^\infty$ seja
uma sequência em $K$.
Então existe $x \in K$ tal que, para todo $\delta > 0$,
$B(x,\delta)$ contém $x_n$ para infinitos $n \in \N$.
Definiremos a subsequência por indução.
A bola $B(x,1)$ contém algum $x_k$; seja $n_1 \coloneqq k$.
Suponha que $n_{j-1}$ esteja definido.
Deve existir $k > n_{j-1}$
tal que $x_k \in B(x,\nicefrac{1}{j})$. Defina
$n_j \coloneqq k$.
Assim, temos a subsequência $\{ x_{n_j} \}_{j=1}^\infty$.
Como
$d(x,x_{n_j}) < \nicefrac{1}{j}$, a \propref{prop:msconvifa} afirma que
$\lim_{j\to\infty} x_{n_j} = x$.

Para a implicação contrária, suponha que toda sequência em $K$
tenha uma subsequência que converge em $K$.
Considere uma cobertura aberta $\{ U_\lambda \}_{\lambda \in I}$ de $K$.
Pelo lema da cobertura de Lebesgue acima, existe $\delta > 0$
tal que, para todo $x \in K$, existe $\lambda \in I$ com
$B(x,\delta) \subset U_\lambda$.

Escolha $x_1 \in K$ e encontre $\lambda_1 \in I$ tal que
$B(x_1,\delta) \subset U_{\lambda_1}$.
Se $K \subset U_{\lambda_1}$, paramos, pois encontramos uma
subcobertura finita.
Caso contrário, existe
$x_2 \in K \setminus U_{\lambda_1}$.
Note que $d(x_2,x_1) \geq \delta$.
Existe algum $\lambda_2 \in I$ tal que
$B(x_2,\delta) \subset U_{\lambda_2}$.
Prosseguimos por indução. Suponha que $\lambda_{n-1}$ esteja definido.
Ou
$U_{\lambda_1} \cup
U_{\lambda_2} \cup \cdots \cup
U_{\lambda_{n-1}}$ é uma cobertura finita de $K$, caso em que
paramos; ou existe
$x_n \in K \setminus \bigl( U_{\lambda_1} \cup
U_{\lambda_2} \cup \cdots \cup
U_{\lambda_{n-1}}\bigr)$.
Note que $d(x_n,x_j) \geq \delta$ para todo $j = 1,2,\ldots,n-1$.
Em seguida, existe algum $\lambda_n \in I$
tal que $B(x_n,\delta) \subset U_{\lambda_n}$.
Veja a \figureref{fig:seqcompactiscompact}.

\begin{myfigureht}
\myincludepdft{seqcompactiscompact}{%
Uma região sombreada com fronteira contínua está identificada por K.
Quatro pontos de K estão identificados como x de índice 1, x de índice 2,
x de índice 3 e x de índice 4.
Círculos de raio delta estão desenhados com centros em x de índice 1,
x de índice 2 e x de índice 3.
Ao redor de cada círculo, há três regiões delimitadas por linhas pontilhadas,
identificadas como U de índice lambda 1,
U de índice lambda 2 e
U de índice lambda 3.
Os discos e, portanto, as regiões ao redor deles parecem cobrir lentamente
todo o conjunto K no lado esquerdo da imagem.}
\caption{Cobertura de $K$ por $U_{\lambda}$. Estão desenhados os pontos
$x_1,x_2,x_3,x_4$,
os três conjuntos
$U_{\lambda_1}$,
$U_{\lambda_2}$,
$U_{\lambda_3}$,
e as três primeiras bolas
de raio $\delta$.\label{fig:seqcompactiscompact}}
\end{myfigureht}

Ou, em algum momento, obtemos uma subcobertura finita de $K$,
ou construímos uma sequência infinita $\{ x_n \}_{n=1}^\infty$ como acima.
Para obter uma contradição, suponha que não haja subcobertura finita e que
tenhamos construído a sequência $\{ x_n \}_{n=1}^\infty$.
Para quaisquer $n$ e $k$ com $n \neq k$,
temos $d(x_n,x_k) \geq \delta$.
Logo, nenhuma subsequência de $\{ x_n \}_{n=1}^\infty$ é de Cauchy.
Portanto, nenhuma subsequência de $\{ x_n \}_{n=1}^\infty$ é convergente,
o que é uma contradição.
\end{proof}

\begin{example} \label{example:intervalcompact}
O \thmref{thm:bwseq},
teorema de Bolzano--Weierstrass para sequências de números reais,
afirma que toda sequência limitada em $\R$ tem uma subsequência convergente.
Portanto, toda sequência em um intervalo fechado $[a,b] \subset \R$ tem
uma subsequência convergente. O limite também pertence a $[a,b]$, pois os
limites preservam desigualdades não estritas. Logo, um intervalo fechado e
limitado $[a,b] \subset \R$ é compacto (sequencialmente).
\end{example}

\begin{prop} \label{prop:closedsubsetcompact}
Seja $(X,d)$ um espaço métrico e seja $K \subset X$ compacto. Se
$E \subset K$ é fechado, então $E$ é compacto.
\end{prop}

Como $K$ é fechado, $E$ é fechado em $K$ se, e somente se, é fechado em $X$.
Veja \propref{prop:topology:subspacesame}.

\begin{proof}
Seja $\{ x_n \}_{n=1}^\infty$ uma sequência em $E$. Ela também é uma
sequência em $K$.
Portanto, tem uma subsequência convergente $\{ x_{n_j} \}_{j=1}^\infty$ cujo
limite é algum $x \in K$. Como $E$ é fechado, o limite de uma sequência em
$E$ também pertence a $E$, logo $x \in E$. Assim, $E$ é compacto.
\end{proof}

\begin{thm}[Heine--Borel%
\footnote{Nomeado em homenagem ao matemático alemão
\href{https://en.wikipedia.org/wiki/Eduard_Heine}{Heinrich Eduard Heine}
(1821--1881)
e ao matemático francês
\href{https://en.wikipedia.org/wiki/\%C3\%89mile_Borel}{F\'elix \'Edouard Justin \'Emile Borel}
(1871--1956).}]%
\index{teorema de Heine--Borel}
\label{thm:msbw}
Todo subconjunto fechado e limitado $K \subset \R^n$ é compacto.
\end{thm}

Assim, os subconjuntos de $\R^n$ são compactos se, e somente se, são
fechados e limitados, uma condição bem mais fácil de verificar.
Reiteramos que o teorema de Heine--Borel vale somente para $\R^n$,
não para espaços métricos em geral.
O teorema não vale sequer para subespaços de $\R^n$, apenas para o próprio
$\R^n$.
Em geral, compacidade implica ser fechado e limitado, mas a recíproca não
vale.

\begin{proof}
Para $\R = \R^1$, suponha que $K \subset \R$ seja fechado e limitado.
Então $K \subset [a,b]$ para algum intervalo fechado e limitado,
que é compacto pelo \exampleref{example:intervalcompact}.
Como $K$ é um subconjunto fechado de um conjunto compacto,
é compacto pela \propref{prop:closedsubsetcompact}.

Faremos a demonstração para $n=2$ e deixaremos o caso de $n$ arbitrário como exercício.
Como $K \subset \R^2$ é limitado, existe um conjunto
$B=[a,b]\times[c,d] \subset \R^2$ tal que $K \subset B$. Mostraremos
que $B$ é compacto. Então $K$, por ser um subconjunto fechado de $B$,
que é compacto, também é compacto.

Seja $\bigl\{ (x_k,y_k) \bigr\}_{k=1}^\infty$ uma sequência em $B$. Isto é,
$a \leq x_k \leq b$ e
$c \leq y_k \leq d$ para todo $k$. Uma sequência limitada de números reais
tem uma subsequência convergente, portanto existe uma subsequência
$\{ x_{k_j} \}_{j=1}^\infty$
convergente. A subsequência
$\{ y_{k_j} \}_{j=1}^\infty$ também é limitada; logo, existe uma
subsequência
$\{ y_{k_{j_i}} \}_{i=1}^\infty$ convergente. Uma subsequência de uma
sequência convergente continua convergente, portanto
$\{ x_{k_{j_i}} \}_{i=1}^\infty$ converge.
Sejam
\begin{equation*}
x \coloneqq \lim_{i\to\infty} x_{k_{j_i}}
\qquad \text{e} \qquad
y \coloneqq \lim_{i\to\infty} y_{k_{j_i}} .
\end{equation*}
Pela \propref{prop:msconveuc},
$\bigl\{ (x_{k_{j_i}},y_{k_{j_i}}) \bigr\}_{i=1}^\infty$ converge a $(x,y)$.
Além disso, como $a \leq x_k \leq b$ e
$c \leq y_k \leq d$ para todo $k$, sabemos que $(x,y) \in B$.
\end{proof}

\begin{example}
A métrica discreta oferece novamente contraexemplos interessantes.
Seja $(X,d)$ um espaço métrico com a métrica discreta, isto é,
$d(x,y) = 1$
se $x \neq y$. Suponha que
$X$ seja infinito. Então
\begin{enumerate}[(i)]
\item $(X,d)$ é um espaço métrico completo.
\item Todo subconjunto $K \subset X$ é fechado e limitado.
\item Um subconjunto $K \subset X$ é compacto se, e somente se, é finito.
\item A conclusão do lema da cobertura de Lebesgue sempre vale,
por exemplo, com $\delta = \nicefrac{1}{2}$,
mesmo para subconjuntos não compactos $K \subset X$.
\end{enumerate}
As demonstrações das afirmações acima ou são triviais, ou ficam como
exercício abaixo.
\end{example}

\begin{remark}
Uma questão sutil sobre sequências de Cauchy, completude, compacidade
e convergência é que a compacidade e a convergência dependem apenas da
topologia, isto é, de quais conjuntos são abertos. Por outro lado,
sequências de Cauchy e completude dependem da métrica propriamente dita.
Veja \exerciseref{exercise:cauchydepndsonmetric}.
\end{remark}

\subsection{Exercícios}

\begin{exercise}
Seja $(X,d)$ um espaço métrico e $A$ um subconjunto finito de $X$.
Mostre que $A$ é compacto.
\end{exercise}

\begin{samepage}
\begin{exercise}
Seja $A \coloneqq \{ \nicefrac{1}{n} : n \in \N \} \subset \R$.
\begin{enumerate}[a)]
\item
Mostre diretamente pela definição que $A$ não é compacto.
\item
Mostre diretamente pela definição que $A \cup \{ 0 \}$ é compacto.
\end{enumerate}
\end{exercise}
\end{samepage}

\begin{exercise}
Seja $(X,d)$ um espaço métrico com a métrica discreta.
\begin{enumerate}[a)]
\item
Demonstre que $X$ é completo.
\item
Demonstre que $X$ é compacto se, e somente se, $X$ é finito.
\end{enumerate}
\end{exercise}

\begin{exercise}
\leavevmode
\begin{enumerate}[a)]
\item
Mostre que a união de um número finito de conjuntos compactos é compacta.
\item
Encontre um exemplo em que a união de infinitos conjuntos compactos não
seja compacta.
\end{enumerate}
\end{exercise}

\begin{exercise}
Demonstre \thmref{thm:msbw} para dimensão arbitrária.
Dica: o segredo é usar a notação correta.
\end{exercise}

\begin{exercise} \label{exercise:subspacecompact}
Mostre que um conjunto compacto $K$ (em qualquer espaço métrico)
é, por si só, um espaço métrico completo (com a métrica de subespaço).
\end{exercise}

\begin{exercise} \label{exercise:CabRcomplete}
Seja $C\bigl([a,b],\R\bigr)$ o espaço métrico definido em
\exampleref{example:msC01}. Mostre que
$C\bigl([a,b],\R\bigr)$ é um espaço métrico completo.
\end{exercise}

\begin{exercise}[Desafio] \label{exercise:msclbounnotcompt}
Seja $C\bigl([0,1],\R\bigr)$ o espaço métrico de
\exampleref{example:msC01}. Seja $0$
a função nula. Mostre que a bola fechada
$C(0,1)$ não é compacta (embora seja fechada e limitada).
Dicas: construa uma sequência de funções contínuas distintas
$\{ f_n \}_{n=1}^\infty$ tal que
$d(f_n,0) = 1$ e $d(f_n,f_k) = 1$ para todo $n \neq k$. Mostre que
o conjunto $\{ f_n : n \in \N \} \subset C(0,1)$ é fechado, mas não compacto.
Busque inspiração em \chapterref{fs:chapter}.
\end{exercise}

\begin{exercise}[Desafio]
Mostre que existe uma métrica em $\R$ que torna $\R$ um conjunto compacto.
\end{exercise}

\begin{exercise}
Suponha que $(X,d)$ seja completo e que tenhamos uma coleção enumerável infinita
de conjuntos compactos não vazios $E_1 \supset E_2 \supset E_3 \supset
\cdots$. Demonstre que $\bigcap_{j=1}^\infty E_j \neq \emptyset$.
\end{exercise}


\begin{exercise}[Desafio]
Seja $C\bigl([0,1],\R\bigr)$ o espaço métrico de \exampleref{example:msC01}.
Seja $K$ o conjunto das funções $f \in C\bigl([0,1],\R\bigr)$ tais que
$f$ é um polinômio quadrático, isto é, $f(x) = a+bx+cx^2$, e tais que
$\babs{f(x)} \leq 1$ para todo $x \in [0,1]$,
ou seja, $f \in C(0,1)$. Mostre que $K$ é compacto.
\end{exercise}

\begin{exercise}[Desafio] \label{exercise:mstotbound}
Seja $(X,d)$ um espaço métrico completo.
Mostre que $K \subset X$ é compacto se, e somente se, $K$ é fechado
e, para todo $\epsilon > 0$,
existe um conjunto finito de pontos $x_1,x_2,\ldots,x_n$ tal que
$K \subset \bigcup_{j=1}^n B(x_j,\epsilon)$.
Observação: Dizemos que um conjunto $K$ assim é \emph{\myindex{totalmente limitado}},
portanto, em um espaço métrico completo, um conjunto é compacto se, e somente
se, é fechado e totalmente limitado.
\end{exercise}

\begin{exercise}
Considere $\N \subset \R$ com a métrica usual. Encontre uma cobertura aberta de $\N$
para a qual a conclusão do lema da cobertura de Lebesgue não vale.
\end{exercise}

\begin{exercise}
Demonstre o \myindex{teorema de Bolzano--Weierstrass} geral:
toda sequência limitada $\{ x_k \}_{k=1}^\infty$ em $\R^n$
tem uma subsequência convergente.
\end{exercise}

\begin{exercise}
Seja $X$ um espaço métrico e seja
$C \subset \sP(X)$ o conjunto dos subconjuntos compactos não vazios de $X$.
Usando a métrica de Hausdorff de \exerciseref{exercise:mshausdorffpseudo},
mostre que $(C,d_H)$ é um espaço métrico. Isto é, mostre que
se $L$ e $K$ são subconjuntos compactos não vazios, então $d_H(L,K) = 0$
se, e somente se, $L=K$.
\end{exercise}

\begin{exercise} \label{exercise:closedcomplete}
Demonstre \propref{prop:closedcomplete}. Isto é,
seja $(X,d)$ um espaço métrico completo e seja $E \subset X$ um conjunto fechado.
Mostre que $E$, com a métrica de subespaço, é um espaço métrico completo.
\end{exercise}

\begin{exercise}
Seja $(X,d)$ um espaço métrico incompleto. Mostre que existe um
conjunto fechado e limitado $E \subset X$ que não é compacto.
\end{exercise}

\begin{exercise}
Seja $(X,d)$ um espaço métrico e $K \subset X$.
Demonstre que $K$ é compacto como subconjunto de $(X,d)$ se, e somente se, $K$ é
compacto como subconjunto de si mesmo com a métrica de subespaço.
\end{exercise}

\begin{exercise} \label{exercise:cauchydepndsonmetric}
Considere duas métricas em $\R$.
Seja $d(x,y) \coloneqq \sabs{x-y}$ a métrica usual,
e seja
$d'(x,y) \coloneqq
\babs{ \frac{x}{1+\sabs{x}} - \frac{y}{1+\sabs{y}} }$.
\begin{enumerate}[a)]
\item
Mostre que $(\R,d')$ é um espaço métrico (se você resolveu
\exerciseref{exercise:extendedrealsmetric}, o cálculo é o mesmo).
\item
Mostre que a topologia é a mesma, isto é, um conjunto é aberto em
$(\R,d)$ se, e somente se, é aberto em $(\R,d')$.
\item
Mostre que um conjunto é compacto em
$(\R,d)$ se, e somente se, é compacto em $(\R,d')$.
\item
Mostre que uma sequência converge em $(\R,d)$ se, e somente se,
converge em $(\R,d')$.
\item
Encontre uma sequência de números reais que seja de Cauchy
em $(\R,d')$, mas não seja de Cauchy em $(\R,d)$.
\item
Embora $(\R,d)$ seja completo, mostre que $(\R,d')$ não é completo.
\end{enumerate}
\end{exercise}

\begin{exercise} \label{exercise:relativelycompactseq}
\pagebreak[2]
Seja $(X,d)$ um espaço métrico completo.
Dizemos que um conjunto $S \subset X$ é \emph{\myindex{relativamente compacto}}
se seu fecho $\widebar{S}$ é compacto.
Demonstre que $S \subset X$ é relativamente compacto se, e somente se,
dada qualquer sequência $\{ x_n \}_{n=1}^\infty$ em $S$, existe uma subsequência
$\{ x_{n_k} \}_{k=1}^\infty$ que converge (em $X$).
\end{exercise}

\sectionnewpage
\section{Funções contínuas}
\label{sec:metcont}

%mbxINTROSUBSECTION

\sectionnotes{1.5--2 lectures}

\subsection{Continuidade}

\begin{defn}
Sejam $(X,d_X)$ e $(Y,d_Y)$ espaços métricos e $c \in X$.
Dizemos que $f \colon X \to Y$ é
\emph{contínua em $c$}\index{contínua em $c$}
se, para todo $\epsilon > 0$,
existe um $\delta > 0$ tal que, sempre que $x \in X$ e $d_X(x,c) <
\delta$, tem-se
$d_Y\bigl(f(x),f(c)\bigr) < \epsilon$.

\medskip

Quando $f \colon X \to Y$ é contínua em todo $c \in X$, dizemos simplesmente
que $f$ é uma \emph{função contínua}\index{função contínua!em um
espaço métrico}\index{função!contínua}.
\end{defn}

A definição coincide com a definição de \chapterref{lim:chapter} quando
$f$ é uma função real definida na reta real — desde que, é claro, usemos
a métrica usual em $\R$.

\begin{prop} \label{prop:contiscont}
Sejam $(X,d_X)$ e $(Y,d_Y)$ espaços métricos.
Então $f \colon X \to Y$ é
contínua em $c \in X$
se, e somente se, para toda sequência $\{ x_n \}_{n=1}^\infty$ em $X$
que converge para $c$, a sequência $\bigl\{ f(x_n) \bigr\}_{n=1}^\infty$ converge
para $f(c)$.
\end{prop}

\begin{proof}
Suponha que $f$ seja contínua em $c$. Seja $\{ x_n \}_{n=1}^\infty$ uma
sequência em $X$ que converge para $c$. Dado $\epsilon > 0$,
existe um $\delta > 0$ tal que $d_X(x,c) < \delta$ implica
$d_Y\bigl(f(x),f(c)\bigr) < \epsilon$. Escolha $M$ de modo que,
para todo $n \geq M$, tenhamos $d_X(x_n,c) < \delta$; então
$d_Y\bigl(f(x_n),f(c)\bigr) < \epsilon$. Portanto, a sequência $\bigl\{ f(x_n) \bigr\}_{n=1}^\infty$
converge para $f(c)$.

Por outro lado, suponha que $f$ não seja contínua em $c$.
Então existe um $\epsilon > 0$
tal que, para todo $n \in \N$, existe um $x_n \in X$
com
$d_X(x_n,c) < \nicefrac{1}{n}$ e tal que $d_Y\bigl(f(x_n),f(c)\bigr) \geq
\epsilon$. Assim, $\{ x_n \}_{n=1}^\infty$ converge para $c$, mas
$\bigl\{ f(x_n) \bigr\}_{n=1}^\infty$
não converge para $f(c)$.
\end{proof}

\begin{example}
Suponha que $f \colon \R^2 \to \R$ seja um polinômio. Isto é,
\begin{equation*}
f(x,y) =
\sum_{j=0}^d
\sum_{k=0}^{d-j}
a_{jk}\,x^jy^k =
a_{0\,0} + a_{1\,0} \, x +
a_{0\,1} \, y+  
a_{2\,0} \, x^2+  
a_{1\,1} \, xy+  
a_{0\,2} \, y^2+ \cdots +
a_{0\,d} \, y^d ,
\end{equation*}
para algum $d \in \N$ (o grau) e $a_{jk} \in \R$. Afirmamos que
$f$ é contínua. Seja $\bigl\{ (x_n,y_n) \bigr\}_{n=1}^\infty$ uma sequência
em $\R^2$ que converge para $(x,y) \in \R^2$. Provamos que isso
significa que $\lim_{n\to\infty} x_n = x$ e $\lim_{n\to\infty} y_n = y$.
Por \propref{prop:contalg},
\begin{equation*}
\lim_{n\to\infty}
f(x_n,y_n) =
\lim_{n\to\infty}
\sum_{j=0}^d
\sum_{k=0}^{d-j}
a_{jk} \, x_n^jy_n^k 
=
\sum_{j=0}^d
\sum_{k=0}^{d-j}
a_{jk} \, x^jy^k
=
f(x,y) .
\end{equation*}
Logo, $f$ é contínua em $(x,y)$; como $(x,y)$ era arbitrário, $f$ é
contínua em toda parte. De modo análogo, um
polinômio em $n$ variáveis é contínuo.
\end{example}

Tenha cuidado ao calcular limites separadamente.
Considere $f \colon \R^2 \to \R$ definida por $f(x,y) \coloneqq \frac{xy}{x^2+y^2}$
fora da origem e por $f(0,0) \coloneqq 0$ na origem.
Veja \figureref{fig:xyxsqysq}.
No \exerciseref{exercise:dicontR2},
pede-se que você prove que $f$ não é contínua na origem.
No entanto, para todo $y$, a função $g(x) \coloneqq f(x,y)$ é
contínua,
e, para todo $x$, a função $h(y) \coloneqq f(x,y)$ é contínua.
\begin{myfigureht}
\myincludegraphics{xyxsqysq}{%
Gráfico de uma superfície que parece comprimida na
origem, de modo que, ao nos aproximarmos da origem no plano xy
em direções diferentes, podemos nos aproximar de valores distintos.}
\caption{Gráfico de $\frac{xy}{x^2+y^2}$.\label{fig:xyxsqysq}}
\end{myfigureht}

\begin{example}
Seja $X$ um espaço métrico e seja $f \colon X \to \C$ uma função com valores
complexos. Escreva $f(p) = g(p) + i \, h(p)$, onde $g \colon X \to \R$
e $h \colon X \to \R$ são as partes real e imaginária de $f$.
Então $f$ é contínua em $c \in X$ se, e somente se, suas partes real
e imaginária são contínuas em~$c$.  
Isso decorre do fato de que a sequência $\bigl\{ f(p_n) = g(p_n) + i \, h(p_n) \bigr\}_{n=1}^\infty$
converge para $f(p) = g(p) + i \, h(p)$ se, e somente se,
$\bigl\{ g(p_n) \bigr\}_{n=1}^\infty$ converge para $g(p)$ e
$\bigl\{ h(p_n) \bigr\}_{n=1}^\infty$ converge para $h(p)$.
\end{example}

\subsection{Compacidade e continuidade}

Aplicações contínuas não levam conjuntos fechados a conjuntos fechados. Por exemplo,
$f \colon (0,1) \to \R$ definida por $f(x) \coloneqq x$ leva o conjunto $(0,1)$, que
é fechado em $(0,1)$, ao conjunto $(0,1)$, que não é fechado em $\R$.
Por outro lado, aplicações contínuas preservam conjuntos compactos.

\begin{lemma} \label{lemma:continuouscompact}
Sejam $(X,d_X)$ e $(Y,d_Y)$ espaços métricos,
e seja $f \colon X \to Y$ uma função contínua. Se
$K \subset X$ é compacto, então $f(K)$ é compacto.
\end{lemma}

\begin{proof}
Uma sequência em $f(K)$ pode ser escrita como
$\bigl\{ f(x_n) \bigr\}_{n=1}^\infty$, onde
$\{ x_n \}_{n=1}^\infty$ é uma sequência em $K$. O conjunto $K$ é compacto e,
portanto, existe uma subsequência
$\{ x_{n_j} \}_{j=1}^\infty$ que converge para algum $x \in K$.
Pela continuidade,
\begin{equation*}
\lim_{j\to\infty} f(x_{n_j}) = f(x) \in f(K) .
\end{equation*}
Assim, toda sequência em $f(K)$ tem uma subsequência que converge
para um ponto de $f(K)$, e $f(K)$ é compacto por \thmref{thm:mscompactisseqcpt}.
\end{proof}

Como antes, dizemos que $f \colon X \to \R$ atinge um
\emph{\myindex{mínimo absoluto}} em $c \in X$ se
\begin{equation*}
f(x) \geq f(c) \qquad \text{para todo } x \in X.
\end{equation*}
Por outro lado, $f$ atinge um
\emph{\myindex{máximo absoluto}} em $c \in X$ se
\begin{equation*}
f(x) \leq f(c) \qquad \text{para todo } x \in X.
\end{equation*}

\begin{thm}
Seja $(X,d)$ um espaço métrico compacto não vazio
e seja $f \colon X \to \R$ contínua. Então
$f$ é limitada e, de fato,
$f$ atinge um mínimo absoluto e um máximo absoluto em $X$.
\end{thm}

\begin{proof}
Como $X$ é compacto e $f$ é contínua,
$f(X) \subset \R$ é compacto. Portanto, $f(X)$ é fechado
e limitado. Em particular,
$\sup f(X) \in f(X)$ e
$\inf f(X) \in f(X)$, pois tanto o supremo quanto o ínfimo
podem ser obtidos por sequências em $f(X)$, e $f(X)$ é fechado.
Portanto, há um $x \in X$ tal que $f(x) = \sup f(X)$
e um $y \in X$ tal que $f(y) = \inf f(X)$.
\end{proof}

\subsection{Continuidade e topologia}

Vejamos como definir a continuidade em termos da topologia, isto é,
dos conjuntos abertos. Já vimos que a topologia determina quais
sequências convergem; portanto, não surpreende que a topologia também
determine a continuidade das funções.

\begin{lemma} \label{lemma:mstopocontloc}
Sejam $(X,d_X)$ e $(Y,d_Y)$ espaços métricos.
Uma função $f \colon X \to Y$ é contínua em $c \in X$
se, e somente se, para toda vizinhança aberta $U$ de $f(c)$ em $Y$, o conjunto
$f^{-1}(U)$ contém uma vizinhança aberta de $c$ em $X$.
Veja \figureref{fig:mscontfuncpt}.
\end{lemma}

\begin{myfigureht}
\myincludepdft{mscontfuncpt}{%
À esquerda, há uma região sombreada escura, marcada com W e com fronteira
pontilhada. Ela está contida em uma região sombreada clara, marcada com a imagem inversa de U e com fronteira traço-ponto. À direita, há uma região
sombreada clara, marcada com U e com fronteira pontilhada. Há um ponto
marcado com c na região W; uma seta marcada com f o liga ao ponto f(c) na região U.}
\caption{Para toda vizinhança $U$ de $f(c)$, o conjunto $f^{-1}(U)$ contém uma vizinhança aberta
$W$ de $c$.\label{fig:mscontfuncpt}}
\end{myfigureht}

\begin{proof}
Suponha primeiro que $f$ seja contínua em $c$.
Seja $U$ uma vizinhança aberta de $f(c)$
em $Y$; então $B_Y\bigl(f(c),\epsilon\bigr) \subset U$ para algum $\epsilon >
0$. Pela continuidade de $f$, existe um $\delta > 0$
tal que, sempre que $x$ satisfaz $d_X(x,c) < \delta$, temos
$d_Y\bigl(f(x),f(c)\bigr) < \epsilon$. Em outras palavras,
\begin{equation*}
B_X(c,\delta) \subset f^{-1}\bigl(B_Y\bigl(f(c),\epsilon\bigr)\bigr) \subset
f^{-1}(U) ,
\end{equation*}
e $B_X(c,\delta)$ é uma vizinhança aberta de $c$.

Para a recíproca,
seja dado $\epsilon > 0$. Se
$f^{-1}\bigl(B_Y\bigl(f(c),\epsilon\bigr)\bigr)$ contém uma vizinhança aberta
$W$ de $c$, então contém uma bola. Isto é, existe algum $\delta > 0$
tal que
\begin{equation*}
B_X(c,\delta) \subset W \subset f^{-1}\bigl(B_Y\bigl(f(c),\epsilon\bigr)\bigr) .
\end{equation*}
Isso significa precisamente que, se $d_X(x,c) < \delta$,
então $d_Y\bigl(f(x),f(c)\bigr) < \epsilon$.
Logo, $f$ é contínua em~$c$.
\end{proof}

\begin{thm} \label{thm:mstopocont}
Sejam $(X,d_X)$ e $(Y,d_Y)$ espaços métricos. Uma função $f \colon X \to Y$
é contínua se, e somente se,
para todo conjunto aberto $U \subset Y$, $f^{-1}(U)$ é aberto em $X$.
\end{thm}

A demonstração decorre de \lemmaref{lemma:mstopocontloc} e fica
como exercício.

\begin{example}
Seja $f \colon X \to Y$ uma função contínua.
O \thmref{thm:mstopocont} nos diz que, se $E \subset Y$ é fechado, então
$f^{-1}(E) = X \setminus f^{-1}(E^c)$ também é fechado. Portanto, se
temos uma função contínua
$f \colon X \to \R$, o
\emph{\myindex{conjunto de zeros}} de $f$, isto é,
$f^{-1}(0) = \bigl\{ x \in X : f(x) = 0 \bigr\}$, é fechado.
Acabamos de demonstrar o resultado mais básico da
\emph{\myindex{geometria algébrica}}, o estudo dos conjuntos
de zeros de polinômios: o conjunto de zeros de um polinômio é fechado.

De modo análogo, o conjunto em que $f$ é não negativa,
$f^{-1}\bigl( [0,\infty) \bigr) = \bigl\{ x \in X : f(x) \geq 0 \bigr\}$,
é fechado. Por outro lado, o
conjunto em que $f$ é positiva,
$f^{-1}\bigl( (0,\infty) \bigr) = \bigl\{ x \in X : f(x) > 0 \bigr\}$,
é aberto.
\end{example}

\subsection{Continuidade uniforme}

Assim como ocorre com funções contínuas na reta real, na definição de continuidade
às vezes é conveniente poder escolher
um único $\delta$ para todos os pontos.

\begin{defn}
Sejam $(X,d_X)$ e $(Y,d_Y)$ espaços métricos.
Dizemos que $f \colon X \to Y$ é
\emph{uniformemente contínua}\index{uniformemente contínua!em um espaço métrico}
se, para todo $\epsilon > 0$,
existe um $\delta > 0$ tal que, sempre que $p,q \in X$ e $d_X(p,q) <
\delta$, temos
$d_Y\bigl(f(p),f(q)\bigr) < \epsilon$.
\end{defn}

Toda função uniformemente contínua é contínua, mas a recíproca não
é necessariamente verdadeira, como vimos.

\begin{thm} \label{thm:Xcompactfunifcont}
Sejam $(X,d_X)$ e $(Y,d_Y)$ espaços métricos.
Suponha que $f \colon X \to Y$ seja contínua e que $X$ seja compacto. Então
$f$ é uniformemente contínua.
\end{thm}

\begin{proof}
Seja dado $\epsilon > 0$. Para cada $c \in X$, escolha $\delta_c > 0$ tal que
$d_Y\bigl(f(x),f(c)\bigr) < \nicefrac{\epsilon}{2}$
sempre que
$x \in B(c,\delta_c)$.
%$d_X(x,c) < \delta_c$.
As bolas
$B(c,\delta_c)$ cobrem $X$, e o espaço $X$ é compacto.
Aplique o \hyperref[ms:lebesgue]{lema da cobertura de Lebesgue} para obter um
$\delta > 0$ tal que, para todo $x \in X$, existe um $c \in X$
para o qual $B(x,\delta) \subset B(c,\delta_c)$.

Suponha que $p, q \in X$ e que $d_X(p,q) < \delta$.
Encontre um $c \in X$ tal que $B(p,\delta) \subset B(c,\delta_c)$.
Então $q \in B(c,\delta_c)$. Pela desigualdade triangular
e pela definição de $\delta_c$,
\begin{equation*}
d_Y\bigl(f(p),f(q)\bigr)
\leq
d_Y\bigl(f(p),f(c)\bigr)
+
d_Y\bigl(f(c),f(q)\bigr)
<
\nicefrac{\epsilon}{2}+
\nicefrac{\epsilon}{2} = \epsilon .  \qedhere
\end{equation*}
\end{proof}

Como aplicação da continuidade uniforme,
demonstraremos um critério útil para a
continuidade de funções definidas por integrais.
Seja $f(x,y)$ uma função de duas variáveis e defina
\begin{equation*}
g(y) \coloneqq \int_a^b f(x,y) \,dx .
\end{equation*}
A questão é: $g$ é contínua?
Estamos perguntando, na verdade, quando duas operações de limite comutam,
o que nem sempre é possível; portanto, é necessária alguma hipótese
adicional. Uma condição suficiente (mas não
necessária) útil é que $f$ seja contínua em um retângulo fechado.

\begin{prop} \label{prop:integralcontcont}
Se $f \colon [a,b] \times [c,d] \to \R$ é contínua,
então $g \colon [c,d] \to \R$, definida por
\begin{equation*}
g(y) \coloneqq \int_a^b f(x,y) \,dx  \qquad \text{é contínua}.
\end{equation*}
\end{prop}

\begin{proof}
Fixe $y \in [c,d]$ e
seja dado $\epsilon > 0$.
Como $f$ é contínua em $[a,b] \times [c,d]$, que é compacto, $f$
é uniformemente contínua.
Em particular, existe um $\delta > 0$ tal que,
sempre que $z \in [c,d]$ e
$\sabs{z-y} < \delta$, temos
$\babs{f(x,z)-f(x,y)} < \frac{\epsilon}{b-a}$ para todo $x \in [a,b]$.
Suponha então que $\sabs{z-y} < \delta$. Então
\begin{multline*}
\babs{
g(z)-
g(y)
}
=
\abs{
\int_a^b 
f(x,z) \,dx 
-
\int_a^b 
f(x,y) \,dx 
}
\\
=
\abs{
\int_a^b 
\bigl(
f(x,z) - f(x,y)
\bigr)
\,dx 
}
\leq
(b-a)
\frac{\epsilon}{b-a}
= \epsilon . \qedhere
\end{multline*}
\end{proof}

Nas aplicações, se estamos interessados na continuidade em $y_0$, basta
aplicar a proposição em $[a,b] \times [y_0-\epsilon,y_0+\epsilon]$
para algum $\epsilon > 0$ pequeno. Por exemplo, se $f$ é contínua em
$[a,b] \times \R$, então $g$ é contínua em $\R$.

\begin{example}
Exemplos úteis de funções uniformemente contínuas são novamente as
\emph{funções Lipschitz}%
\index{função Lipschitz!em um espaço métrico}%
\index{função!Lipschitz}
Isto é, se
$(X,d_X)$ e $(Y,d_Y)$ são espaços métricos, então $f \colon X \to Y$
é chamada função Lipschitz ou $K$-Lipschitz se existe um $K \in \R$ tal que
\begin{equation*}
d_Y\bigl(f(p),f(q)\bigr) \leq K \, d_X(p,q)
\qquad \text{para todos } p,q \in X.
\end{equation*}
Uma função Lipschitz é uniformemente contínua:
basta tomar $\delta = \nicefrac{\epsilon}{K}$ (supondo que $K > 0$).
Uma função pode ser uniformemente contínua
sem ser Lipschitz,
como já vimos: $\sqrt{x}$ em $[0,1]$
é uniformemente contínua, mas não é Lipschitz.

Vale mencionar que,
se uma função é Lipschitz, em geral é
mais simples demonstrar diretamente que ela é Lipschitz, mesmo quando
o único objetivo é saber que ela é contínua.
\end{example}

\subsection{Pontos de acumulação e limites de funções}

Embora ainda não tenhamos iniciado a discussão da continuidade com eles,
e só venhamos a precisar deles
\volIIref{volume II\@,}{bem mais adiante,}
vamos também
traduzir a ideia de limite de uma função na reta real para espaços métricos.
Novamente, precisamos começar pelos pontos de acumulação.

\begin{defn}
Seja $(X,d)$ um espaço métrico e
$S \subset X$. Dizemos que um ponto $p \in X$ é
um \emph{ponto de acumulação}\index{ponto de acumulação!em um espaço métrico} de $S$
se, para todo $\epsilon > 0$, o conjunto $B(p,\epsilon) \cap S
\setminus \{ p \}$ não é vazio.
\end{defn}

Não basta que $p$ pertença ao fecho de $S$;
ele deve pertencer ao fecho de
$S \setminus \{ p \}$ (exercício).
Assim, $p$ é um ponto de acumulação se, e somente se, existe uma sequência em $S \setminus
\{ p \}$ que converge para $p$.

\begin{defn}
\index{limite!de uma função em um espaço métrico}%
Sejam $(X,d_X)$, $(Y,d_Y)$ espaços métricos, $S \subset X$, $p \in X$ um ponto de acumulação de $S$,
e seja $f \colon S \to Y$ uma função.
Suponha que exista um $L \in Y$ e que, para todo $\epsilon > 0$,
exista um $\delta > 0$ tal que, sempre que $x \in S \setminus \{ p \}$
e $d_X(x,p) < \delta$, temos
\begin{equation*}
d_Y\bigl(f(x),L\bigr) < \epsilon .
\end{equation*}
Então dizemos que $f(x)$
\emph{converge}\index{converge!função em um espaço métrico} para $L$ quando $x$ tende
a $p$, e que $L$ é um \emph{limite} de $f(x)$ quando $x$
tende a $p$. Se $L$ é único, escrevemos
\glsadd{not:limfunc}
\begin{equation*}
\lim_{x \to p} f(x) \coloneqq L .
\end{equation*}
Se $f(x)$ não converge quando $x$ tende a $p$, dizemos que $f$
\emph{diverge}\index{diverge!função em um espaço métrico} em $p$.
\end{defn}

Como de costume, provamos que o limite, se existir, é único.
A demonstração é uma tradução direta da demonstração
de \chapterref{lim:chapter}; por isso, deixamos o resultado como exercício.

\begin{prop} \label{prop:mslimitisunique}
Sejam $(X,d_X)$ e $(Y,d_Y)$ espaços métricos, $S \subset X$, $p \in X$
um ponto de acumulação de $S$, e seja $f \colon S \to Y$ uma função
tal que $f(x)$ converge quando $x$ tende a $p$. Então
o limite de $f(x)$ quando $x$ tende a $p$ é único.
\end{prop}

Em qualquer espaço métrico, assim como em $\R$, limites de funções podem ser
substituídos por limites sequenciais. A demonstração também é uma tradução direta
da demonstração de \chapterref{lim:chapter}, e a deixamos como
exercício. Em resumo, precisamos demonstrar os resultados apenas para
limites sequenciais.

\begin{lemma}\label{ms:seqflimit:lemma}
Sejam $(X,d_X)$ e $(Y,d_Y)$ espaços métricos, $S \subset X$, $p \in X$
um ponto de acumulação de $S$, e seja $f \colon S \to Y$ uma função.

Então
$f(x)$ converge para $L \in Y$ quando $x$ tende a $p$ se, e somente se, para toda
sequência $\{ x_n \}_{n=1}^\infty$
em $S \setminus \{p\}$
tal que $\lim_{n\to\infty} x_n = p$,
a sequência $\bigl\{ f(x_n) \bigr\}_{n=1}^\infty$ converge para $L$.
\end{lemma}

Aplicando \propref{prop:contiscont} ou diretamente a definição, concluímos
(exercício), como em \chapterref{lim:chapter}, que, para pontos de acumulação $p$ de $S
\subset X$, a função
$f \colon S \to Y$ é contínua em $p$ se, e somente se,
\begin{equation*}
\lim_{x \to p} f(x) = f(p) .
\end{equation*}

\subsection{Exercícios}

\begin{exercise}
Considere $\N \subset \R$ com a métrica usual. Seja $(X,d)$ um
espaço métrico e seja $f \colon X \to \N$ uma função contínua.
\begin{enumerate}[a)]
\item
Demonstre que, se $X$ é conexo,
então $f$ é constante (a imagem de $f$ tem um único elemento).
\item
Dê um exemplo em que $X$ é desconexo e $f$ não é constante.
\end{enumerate}
\end{exercise}

\begin{exercise} \label{exercise:dicontR2}
\pagebreak[2]
Defina $f \colon \R^2 \to \R$ por $f(0,0) \coloneqq 0$ e
$f(x,y) \coloneqq \frac{xy}{x^2+y^2}$ se $(x,y) \neq (0,0)$.
Veja \figureref{fig:xyxsqysq}.
\begin{enumerate}[a)]
\item
Demonstre que, para todo $x$ fixado,
a função que leva $y$ a $f(x,y)$ é contínua. De modo análogo,
para todo $y$ fixado, a função que leva $x$ a $f(x,y)$ é contínua.
\item
Demonstre que $f$ não é contínua.
\end{enumerate}
\end{exercise}

\begin{exercise}
\pagebreak[2]
Suponha que $(X,d_X)$ e $(Y,d_Y)$ sejam espaços métricos e que
$f \colon X \to Y$ seja contínua.
Seja $A \subset X$.
\begin{enumerate}[a)]
\item
Demonstre que $f(\widebar{A}) \subset \overline{f(A)}$.
\item
Dê um exemplo em que $f(\widebar{A})$ seja um subconjunto próprio de
$\overline{f(A)}$.
\end{enumerate}
\end{exercise}

\begin{exercise}
Demonstre o \thmref{thm:mstopocont}. Dica: use \lemmaref{lemma:mstopocontloc}.
\end{exercise}

\begin{exercise} \label{exercise:msconnconn}
Suponha que $f \colon X \to Y$ seja contínua para os espaços métricos $(X,d_X)$
e $(Y,d_Y)$. Demonstre que, se $X$ é conexo, então $f(X)$ é conexo.
\end{exercise}

\begin{exercise}
Demonstre a seguinte versão do
\hyperref[IVT:thm]{teorema do valor intermediário}. Seja $(X,d)$ um espaço métrico
conexo e seja $f \colon X \to \R$ uma função contínua. Suponha que
$x_0,x_1 \in X$ e $y \in \R$ sejam tais que $f(x_0) < y < f(x_1)$.
Demonstre então que existe um $z \in X$ tal que $f(z) = y$.
Dica: veja \exerciseref{exercise:msconnconn}.
\end{exercise}

\begin{exercise}
Uma função contínua $f \colon X \to Y$ entre os espaços métricos $(X,d_X)$ e
$(Y,d_Y)$ é dita \emph{\myindex{própria}}
se, para todo conjunto compacto $K \subset Y$, o conjunto $f^{-1}(K)$ é compacto.
Suponha que uma função contínua $f \colon (0,1) \to (0,1)$ seja própria e que
$\{ x_n \}_{n=1}^\infty$ seja uma sequência em $(0,1)$ que converge para $0$. Demonstre que
$\bigl\{ f(x_n) \bigr\}_{n=1}^\infty$ não tem subsequência que
convirja em $(0,1)$.
\end{exercise}

\begin{exercise}
Sejam $(X,d_X)$ e $(Y,d_Y)$ espaços métricos e seja
$f \colon X \to Y$ uma função contínua, bijetora. Suponha
que $X$ seja compacto. Demonstre que a inversa $f^{-1} \colon Y \to X$
é contínua.
\end{exercise}

\begin{exercise}
Considere o espaço métrico das funções contínuas $C\bigl([0,1],\R\bigr)$.
Seja
$k \colon [0,1] \times [0,1] \to \R$ uma função contínua.
Dada uma função $f \in C\bigl([0,1],\R\bigr)$, defina
\begin{equation*}
\varphi_f(x) \coloneqq \int_0^1 k(x,y) f(y) \,dy .
\end{equation*}
\begin{enumerate}[a)]
\item
Demonstre que $T(f) \coloneqq \varphi_f$ define uma função
$T \colon C\bigl([0,1],\R\bigr) \to C\bigl([0,1],\R\bigr)$.
\item
Demonstre que $T$ é contínua.
\end{enumerate}
\end{exercise}

\begin{samepage}
\begin{exercise}
Seja $(X,d)$ um espaço métrico.
\begin{enumerate}[a)]
\item
Se $p \in X$,
demonstre que $f \colon X \to \R$, definida
por $f(x) \coloneqq d(x,p)$, é contínua.
\item
Defina uma métrica em $X \times X$ como em \exerciseref{exercise:mscross}, item
b, e demonstre que $g \colon X \times X \to \R$, definida por
$g(x,y) \coloneqq d(x,y)$, é contínua.
\item
Demonstre que, se $K_1$ e $K_2$ são subconjuntos compactos de $X$, então
existem $p \in K_1$ e $q \in K_2$ tais que $d(p,q)$ é mínimo,
isto é, $d(p,q) = \inf \{ d(x,y) \colon x \in K_1, y \in K_2 \}$.
\end{enumerate}
\end{exercise}
\end{samepage}

\begin{exercise}
Seja $(X,d)$ um espaço métrico compacto e seja $C(X,\R)$ o conjunto
das funções contínuas com valores reais. Defina
\begin{equation*}
d(f,g) \coloneqq \snorm{f-g}_X = \sup_{x \in X} \babs{f(x)-g(x)} .
\end{equation*}
\begin{enumerate}[a)]
\item
Demonstre que $d$ torna $C(X,\R)$ um espaço métrico.
\item
Demonstre que, para todo $x \in X$, a função avaliação
$E_x \colon C(X,\R) \to \R$, definida por $E_x(f) \coloneqq f(x)$,
é contínua.
\end{enumerate}
\end{exercise}

\begin{samepage}
\begin{exercise}
Seja $C\bigl([a,b],\R\bigr)$ o conjunto das funções contínuas e
$C^1\bigl([a,b],\R\bigr)$\glsadd{not:contdifffuncs} o conjunto das
funções continuamente diferenciáveis em $[a,b]$.
Defina
\begin{equation*}
d_{C}(f,g) \coloneqq \snorm{f-g}_{[a,b]}
\qquad \text{e} \qquad
d_{C^1}(f,g) \coloneqq \snorm{f-g}_{[a,b]} + \snorm{f'-g'}_{[a,b]},
\end{equation*}
onde $\snorm{\cdot}_{[a,b]}$ é a norma uniforme.
Por \exampleref{example:msC01} e \exerciseref{exercise:C1ab}, sabemos que
$C\bigl([a,b],\R\bigr)$ com $d_C$ é um espaço métrico, e
também o é
$C^1\bigl([a,b],\R\bigr)$ com $d_{C^1}$.
\begin{enumerate}[a)]
\item
Demonstre que o operador derivada $D \colon 
C^1\bigl([a,b],\R\bigr) \to C\bigl([a,b],\R\bigr)$, definido por
$D(f) \coloneqq f'$, é contínuo.
\item
Por outro lado, se considerarmos a métrica $d_C$ em $C^1\bigl([a,b],\R\bigr)$,
demonstre que o operador derivada deixa de ser contínuo. Dica: considere
$\sin(n x)$.
\end{enumerate}
\end{exercise}
\end{samepage}

\begin{exercise}
Sejam $(X,d)$ um espaço métrico, $S \subset X$ e $p \in X$. Demonstre que
$p$ é um ponto de acumulação de $S$ se, e somente se, $p \in \overline{S \setminus \{
p \}}$.
\end{exercise}

\begin{exercise}
Demonstre \propref{prop:mslimitisunique}.
\end{exercise}

\begin{exercise}
Demonstre \lemmaref{ms:seqflimit:lemma}.
\end{exercise}

\begin{exercise}
Sejam $(X,d_X)$ e $(Y,d_Y)$ espaços métricos, $S \subset X$, $p \in X$
um ponto de acumulação de $S$, e seja $f \colon S \to Y$ uma função.
Demonstre que
$f \colon S \to Y$ é contínua em $p$ se, e somente se,
\begin{equation*}
\lim_{x \to p} f(x) = f(p) .
\end{equation*}
\end{exercise}

\begin{exercise}
Defina
\begin{equation*}
f(x,y) \coloneqq
\begin{cases}
\frac{2xy}{x^4+y^2} & \text{se } (x,y) \neq (0,0), \\
0 & \text{se } (x,y) = (0,0) .
\end{cases}
\end{equation*}
\begin{enumerate}[a)]
\item
Demonstre que, para todo $y$ fixado, a função que leva $x$ a $f(x,y)$
é contínua e, portanto, integrável no sentido de Riemann.
\item
Para todo $x$ fixado, a função que leva $y$ a $f(x,y)$ é contínua.
\item
Demonstre que $f$ não é contínua em $(0,0)$.
\item
Agora demonstre que $g(y) \coloneqq \int_0^1 f(x,y)\,dx$ não é contínua em $y=0$.
\end{enumerate}
Observação: fique à vontade para usar o que sabe sobre $\arctan$ do cálculo,
em particular que $\frac{d}{ds} \bigl[ \arctan(s) \bigr] = \frac{1}{1+s^2}$.
\end{exercise}

\begin{exercise} \label{exercise:integralcontcontextra}
Demonstre uma versão mais forte da \propref{prop:integralcontcont}:
se $f \colon (a,b) \times (c,d) \to \R$ é uma função contínua limitada,
então $g \colon (c,d) \to \R$, definida por
\begin{equation*}
g(y) \coloneqq \int_a^b f(x,y) \,dx  \qquad \text{é contínua}.
\end{equation*}
Dica: primeiro integre em $[a+\nicefrac{1}{n},b-\nicefrac{1}{n}]$.
\end{exercise}

%%%%%%%%%%%%%%%%%%%%%%%%%%%%%%%%%%%%%%%%%%%%%%%%%%%%%%%%%%%%%%%%%%%%%%%%%%%%%%

\sectionnewpage
\section{Fixed point theorem and Picard's theorem again}
\label{sec:metpicard}

%mbxINTROSUBSECTION

\sectionnotes{1 lecture (optional, does not require \sectionref{sec:picard})}

In this section we prove the fixed point theorem for contraction
mappings.  As an application we prove Picard's theorem, which we proved
without metric spaces in \sectionref{sec:picard}.
The proof presented here is similar, but it goes much
more smoothly with metric spaces and the fixed point theorem.

\subsection{Fixed point theorem}

\begin{defn}
Let $(X,d_X)$ and $(Y,d_Y)$ be metric spaces.
A map
$\varphi \colon X \to Y$ is a \emph{\myindex{contraction}}
(or a contractive map) if it is
a $k$-Lipschitz map for some $k < 1$, i.e.\ if there exists a $k < 1$ such that
\begin{equation*}
d_Y\bigl(\varphi(p),\varphi(q)\bigr) \leq k\, d_X(p,q)
\qquad \text{for all } p,q \in X.
\end{equation*}

\medskip

Given a map $\varphi \colon X \to X$, a point $x \in X$ is called a
\emph{\myindex{fixed point}}
if $\varphi(x)=x$.
\end{defn}

\begin{thm}%
[Contraction mapping principle\index{contraction mapping principle}
or \myindex{Banach fixed point theorem}\index{fixed point theorem}%
\footnote{Named after the Polish mathematician
\href{https://en.wikipedia.org/wiki/Stefan_Banach}{Stefan Banach}
(1892--1945) who first stated the theorem in 1922.}]
\label{thm:contr}
Let $(X,d)$ be a nonempty complete metric space and $\varphi \colon X \to X$ a
contraction.
Then $\varphi$ has a unique fixed point.
\end{thm}

The words \emph{complete} and \emph{contraction} are necessary.
See \exerciseref{exercise:nofixedpoint}.

\begin{proof}
Pick $x_0 \in X$.
Define a sequence $\{ x_n \}_{n=1}^\infty$ by $x_{n+1} \coloneqq \varphi(x_n)$.  Then
\begin{equation*}
d(x_{n+1},x_n) = d\bigl(\varphi(x_n),\varphi(x_{n-1})\bigr)
\leq k d(x_n,x_{n-1}) .
%\leq \cdots
%\leq k^n d(x_1,x_0) .
\end{equation*}
Repeating $n$ times, we get $d(x_{n+1},x_n) \leq k^n d(x_1,x_0)$.
For $m > n$,
\begin{equation*}
\begin{split}
d(x_m,x_n)
& \leq \sum_{i=n}^{m-1} d(x_{i+1},x_i) \\
& \leq \sum_{i=n}^{m-1} k^i d(x_1,x_0) \\
& = k^n d(x_1,x_0) \sum_{i=0}^{m-n-1} k^i \\
& \leq k^n d(x_1,x_0) \sum_{i=0}^{\infty} k^i
= k^n d(x_1,x_0) \frac{1}{1-k} .
\end{split}
\end{equation*}
In particular, the sequence is Cauchy (why?).  Since $X$ is complete,
we let $x \coloneqq \lim_{n\to\infty} x_n$, and we claim that $x$
is our unique fixed point.

Fixed point?  The function $\varphi$ is a contraction,
so it is Lipschitz continuous:
\begin{equation*}
\varphi(x) = \varphi\Bigl( \lim_{n\to\infty} x_n\Bigr) = \lim_{n\to\infty}
\varphi(x_n) =
\lim_{n\to\infty} x_{n+1} = x .
\end{equation*}
Unique?  Let $x$ and $y$ be fixed points.
\begin{equation*}
d(x,y) = d\bigl(\varphi(x),\varphi(y)\bigr) \leq k\, d(x,y) .
\end{equation*}
As $k < 1$, the inequality means that $d(x,y) = 0$, and hence $x=y$.  The theorem is
proved.
\end{proof}

The proof is constructive.  Not only do we know 
a unique fixed point exists, we know how to find it.  Start with
any point $x_0 \in X$, then iterate $\varphi(x_0)$,
$\varphi(\varphi(x_0))$,
$\varphi(\varphi(\varphi(x_0)))$, etc.\ to find better and better approximations.
We can even find how far away
from the fixed point we are.
See the exercises.  The idea of the proof is
therefore useful in real-world applications.

\subsection{Picard's theorem}

We start with the metric space where we will apply the
fixed point theorem:
the space $C\bigl([a,b],\R\bigr)$ of \exampleref{example:msC01},
the space of continuous functions $f \colon [a,b] \to \R$ with the metric
\begin{equation*}
d(f,g) \coloneqq \snorm{f-g}_{[a,b]} = \sup_{x \in [a,b]} \babs{f(x)-g(x)} .
\end{equation*}
Convergence in this metric is convergence in uniform norm, or in other
words, uniform convergence.  Therefore,
$C\bigl([a,b],\R\bigr)$ is a complete metric space,
see \propref{prop:CabRcomplete}.

\medskip

Consider now the ordinary differential equation
\begin{equation*}
\frac{dy}{dx} = F(x,y) .
\end{equation*}
Given some $x_0, y_0$, we desire a function $y=f(x)$ such that
$f(x_0) = y_0$ and such that
\begin{equation*}
f'(x) = F\bigl(x,f(x)\bigr) .
\end{equation*}
To avoid having to come up with many names, we often simply write $y' = F(x,y)$
for the equation
and $y(x)$ for the solution.

The simplest example is the equation $y' = y$, $y(0) = 1$.
The solution is the exponential $y(x) = e^x$.  A somewhat more complicated
example is $y' = -2xy$, $y(0) = 1$, whose solution is the Gaussian
$y(x) = e^{-x^2}$.

A subtle issue is how long the solution exists.
Consider the equation $y' = y^2$, $y(0)=1$.  Then $y(x) = \frac{1}{1-x}$ is a
solution.  While $F$ is a reasonably \myquote{nice} function and in particular
it exists for all $x$ and $y$, the solution \myquote{blows up} at $x=1$.
For more examples related to Picard's theorem, see \sectionref{sec:picard}.

We will look for the solution in 
$C\bigl([a,b],\R\bigr)$, which may feel strange at first as
we are searching for a differentiable function.
The explanation is that we consider the
corresponding integral equation
\begin{equation*}
f(x)
=
y_0 + \int_{x_0}^x F\bigl(t,f(t)\bigr)\,dt .
\end{equation*}
To solve this integral equation, we only need a continuous function, so
in some sense our task should be easier---we have more candidate functions
to try.  This way of thinking is quite typical when solving differential
equations.

\begin{samepage}
\begin{thm}[Picard's theorem on existence and uniqueness]%
\index{existence and uniqueness theorem}\index{Picard's theorem}
Let $I, J \subset \R$ be closed and bounded intervals,
let $I^\circ$ and $J^\circ$ be their interiors, and 
let $(x_0,y_0) \in I^\circ \times J^\circ$.
Suppose $F \colon I \times J \to \R$ is continuous
and Lipschitz in the second variable, that is, there exists
an $L \in \R$ such that
\begin{equation*}
\babs{F(x,y) - F(x,z)} \leq L \sabs{y-z}
\qquad \text{for all } y,z \in J, x \in I .
\end{equation*}
Then there exists an $h > 0$ such that $[x_0-h,x_0+h] \subset I$ and a unique differentiable
function $f \colon [x_0 - h, x_0 + h] \to J \subset \R$ such that
\begin{equation*}
f'(x) = F\bigl(x,f(x)\bigr) \qquad \text{and} \qquad f(x_0) = y_0.
\end{equation*}
\end{thm}
\end{samepage}

\begin{proof}
Without loss of generality, assume $x_0 =0$ (exercise).
As $I \times J$ is compact and
$F$ is continuous, $F$ is bounded.
So find an $M > 0$ such that
$\babs{F(x,y)} \leq M$ for all $(x,y) \in I\times J$.
Pick $\alpha > 0$ such that
$[-\alpha,\alpha] \subset I$ and $[y_0-\alpha, y_0 + \alpha] \subset J$.
Let
\begin{equation*}
h \coloneqq \min \left\{ \alpha, \frac{\alpha}{M+L\alpha} \right\} .
\end{equation*}
Note $[-h,h] \subset I$.  Let
\begin{equation*}
Y \coloneqq \Bigl\{ f \in C\bigl([-h,h],\R\bigr) : f\bigl([-h,h]\bigr) \subset J \Bigr\} .
\end{equation*}
That is, $Y$ is the set of continuous functions on $[-h,h]$ with values in
$J$, in other words,
exactly those functions where $F\bigl(x,f(x)\bigr)$ makes sense.
It is left as an exercise to show that $Y$ is a closed subset of $C\bigl([-h,h],\R\bigr)$
(because $J$ is closed).
The space $C\bigl([-h,h],\R\bigr)$ is complete, and
a closed subset of a complete metric space is a complete metric space with
the subspace metric, see \propref{prop:closedcomplete}.  So $Y$ with the
subspace metric is a complete metric space.
We will write $d(f,g) = \snorm{f-g}_{[-h,h]}$ for this metric.

Define a mapping
$T \colon Y \to C\bigl([-h,h],\R\bigr)$ by
\begin{equation*}
T(f)(x)
\coloneqq
y_0 + \int_0^x F\bigl(t,f(t)\bigr)\,dt .
\end{equation*}
It is an exercise to check that
$T$ is well-defined, and that for $f \in Y$, $T(f)$ really is in $C\bigl([-h,h],\R\bigr)$.
Let $f \in Y$ and $\sabs{x} \leq h$.
As $F$ is bounded by $M$, we have
\begin{equation*}
\begin{split}
\babs{T(f)(x) - y_0}
&= \abs{\int_0^x F\bigl(t,f(t)\bigr)\,dt} \\
& \leq 
\sabs{x} M \leq hM \leq \frac{\alpha M}{M+ L\alpha} \leq \alpha .
\end{split}
\end{equation*}
So $T(f)\bigl([-h,h]\bigr) \subset [y_0-\alpha,y_0+\alpha] \subset J$, and
$T(f) \in Y$.  In other words, $T(Y) \subset Y$.  From now on,
we consider $T$ as a mapping of $Y$ to $Y$.

We claim $T \colon Y \to Y$ is a contraction.  First, for $x \in [-h,h]$
and $f,g \in Y$, we have
\begin{equation*}
\Babs{F\bigl(x,f(x)\bigr) - F\bigl(x,g(x)\bigr)} \leq
L\babs{f(x)- g(x)} \leq L \, d(f,g) .
\end{equation*}
Therefore,
\begin{equation*}
\begin{split}
\babs{T(f)(x) - T(g)(x)}
&= \abs{\int_0^x \Bigl( F\bigl(t,f(t)\bigr) - F\bigl(t,g(t)\bigr)\Bigr)\,dt} \\
& \leq \sabs{x} L \, d(f,g)
 \leq h L\, d(f,g)
 \leq \frac{L\alpha}{M+L\alpha} \, d(f,g) .
\end{split}
\end{equation*}
We chose $M > 0$ and so
$\frac{L\alpha}{M+L\alpha} < 1$.
Take supremum over $x \in [-h,h]$ of the left-hand side above to obtain
$d\bigl(T(f),T(g)\bigr) \leq \frac{L\alpha}{M+L\alpha} \, d(f,g)$,
that is, $T$ is a contraction.

The fixed point theorem (\thmref{thm:contr})
gives a unique $f \in Y$ such that $T(f) = f$.
In other words,
\begin{equation*} %\label{equation:msinteqpicard}
f(x) = y_0 + \int_0^x F\bigl(t,f(t)\bigr)\,dt .
\end{equation*}
Clearly, $f(0) = y_0$.
By the fundamental theorem of calculus (\thmref{thm:FTCv2}),
$f$ is differentiable and its derivative is
$F\bigl(x,f(x)\bigr)$.
Differentiable functions are continuous, so
$f$ is the unique differentiable $f \colon [-h,h] \to J$
such that
 $f'(x) = F\bigl(x,f(x)\bigr)$ and $f(0) = y_0$.
\end{proof}

\subsection{Exercises}

\begin{exnote}
For more exercises related to Picard's theorem see \sectionref{sec:picard}.
\end{exnote}

\begin{exercise}
Let $J$ be a closed and bounded interval and
$Y \coloneqq \bigl\{ f \in C\bigl([-h,h],\R\bigr) : f\bigl([-h,h]\bigr) \subset J \bigr\}$.
Show that $Y \subset C\bigl([-h,h],\R\bigr)$ is closed.  Hint: $J$ is closed.
\end{exercise}

\begin{exercise}
In the proof of Picard's theorem,
show that if $f \colon [-h,h] \to J$ is continuous, then $F\bigl(t,f(t)\bigr)$
is continuous on $[-h,h]$ as a function of $t$.  Use this to show that
\begin{equation*}
T(f)(x)
\coloneqq
y_0 + \int_0^x F\bigl(t,f(t)\bigr)\,dt
\end{equation*}
is well-defined and that $T(f) \in C\bigl([-h,h],\R\bigr)$.
\end{exercise}


\begin{exercise}
Prove that in the proof of Picard's theorem,
the statement \myquote{Without loss of generality assume $x_0 = 0$} is
justified.  That is, prove that if we know the theorem with $x_0 = 0$, the
theorem is true as stated.
\end{exercise}

\begin{exercise}
\pagebreak[2]
Let $F \colon \R \to \R$ be defined by
$F(x) \coloneqq kx + b$ where $0 < k < 1$, $b \in \R$.
\begin{enumerate}[a)]
\item
Show that $F$ is a contraction.
\item
Find the fixed point and show directly that it is unique.
\end{enumerate}
\end{exercise}

\begin{exercise}
\pagebreak[2]
Let $f \colon [0,\nicefrac{1}{4}] \to [0,\nicefrac{1}{4}]$ be defined by
$f(x) \coloneqq x^2$.
\begin{enumerate}[a)]
\item
Show that $f$
is a contraction, and find the best (smallest) $k$ from the definition that works.
\item
Find the fixed point and show directly that it is unique.
\end{enumerate}
\end{exercise}

\begin{samepage}
\begin{exercise} \label{exercise:nofixedpoint}
\leavevmode
\begin{enumerate}[a)]
\item
Find an example of a contraction $f \colon X \to X$
of a non-complete metric space $X$ with no
fixed point.
\item
Find a 1-Lipschitz map $f \colon X \to X$ of a complete metric space $X$ with no fixed point.
\end{enumerate}
\end{exercise}
\end{samepage}

\begin{exercise}
Consider $y' =y^2$, $y(0)=1$.  Use the iteration scheme
from the proof of the contraction mapping principle.
Start with $f_0(x) = 1$.  Find a 
few iterates (at least up to $f_2$).  Prove that
the pointwise limit of $f_n$ is $\frac{1}{1-x}$, that is, for every $x$
with $\sabs{x} < h$ for some $h > 0$,
prove that $\lim\limits_{n\to\infty}f_n(x) = \frac{1}{1-x}$.
\end{exercise}

\begin{exercise}
Suppose $f \colon X \to X$ is a contraction for $k < 1$.  Suppose you use the iteration
procedure with $x_{n+1} \coloneqq f(x_n)$ as in the proof of the fixed point theorem.
Suppose $x$ is the fixed
point of $f$.
\begin{enumerate}[a)]
\item
Show that $d(x,x_n) \leq k^n d(x_1,x_0) \frac{1}{1-k}$ for all $n \in \N$.
\item
Suppose $d(y_1,y_2) \leq 16$ for all $y_1,y_2 \in X$, and $k=
\nicefrac{1}{2}$.  Find an $N$ such that starting at any given point $x_0 \in X$, 
$d(x,x_n) \leq 2^{-16}$ for all $n \geq N$.
\end{enumerate}
\end{exercise}

\begin{exercise}
Let $f(x) \coloneqq x-\frac{x^2-2}{2x}$ (you may recognize Newton's method for
$\sqrt{2}$).
\begin{enumerate}[a)]
\item
Prove $f\bigl([1,\infty)\bigr) \subset [1,\infty)$.
\item
Prove that $f \colon [1,\infty) \to [1,\infty)$ is a contraction.
\item
Show that the fixed point theorem applies,
find the unique $x \geq 1$ such that $f(x) = x$,
and show that $x = \sqrt{2}$.
Note: In particular, the technique from the proof of the theorem
can be used to approximate $\sqrt{2}$.
\end{enumerate}
\end{exercise}

\begin{exercise}
Suppose $f \colon X \to X$ is a contraction, and $(X,d)$ is a metric space
with the discrete metric, that is, $d(x,y) = 1$ whenever $x \neq y$.
Show that $f$ is constant, that is,
there exists a $c \in X$ such that $f(x) = c$ for all $x \in X$.
\end{exercise}

\begin{exercise}
Suppose $(X,d)$ is a nonempty complete metric space,
$f \colon X \to X$ is a mapping, and denote
by $f^n$ the $n$th iterate of $f$.  Suppose
for every $n$ there exists a $k_n > 0$ such that
$d\bigl(f^n(x),f^n(y)\bigr) \leq k_n \, d(x,y)$
for all $x,y \in X$,
where $\sum_{n=1}^\infty k_n < \infty$.
Prove that $f$ has a unique fixed point in $X$.
\end{exercise}
